% !TeX program = pdflatex
% Unified manuscript: all proofs, appendices, and references are embedded.
% Compile with pdflatex until references stabilize (usually two or three passes).
% No external bibliography or figure files are required.

\documentclass[11pt]{article}
\usepackage[T1]{fontenc}
\usepackage[utf8]{inputenc}
\usepackage{lmodern}
\usepackage[a4paper,margin=1in]{geometry}
\usepackage{amsmath,amssymb,amsthm,mathtools}
\usepackage{aliascnt}
\usepackage{microtype}
\usepackage{booktabs,array,tabularx}
\usepackage[font=small,labelfont=bf]{caption}
\usepackage{enumitem}
\usepackage[numbers,sort&compress]{natbib}
\usepackage[hidelinks]{hyperref}
\usepackage[nameinlink,noabbrev]{cleveref}
\usepackage{fancyhdr}
\usepackage{needspace}
\hypersetup{pdftitle={Near-Optimal Complexity of Convex Optimization with Exact Coordinate Subgradient Oracles},pdfauthor={Colombo Team},pdfsubject={Convex optimization with exact coordinate subgradients},pdfkeywords={convex optimization, coordinate subgradient, oracle complexity, randomized lower bounds}}
\setlength{\headheight}{14pt}
\pagestyle{fancy}
\fancyhf{}
\fancyhead[L]{\small Exact coordinate subgradient oracles}
\fancyhead[R]{\small Convex Lipschitz optimization}
\fancyfoot[C]{\thepage}
\renewcommand{\headrulewidth}{0.3pt}
\setlength{\emergencystretch}{2em}
\newcolumntype{Y}{>{\raggedright\arraybackslash}X}
\newcolumntype{L}[1]{>{\raggedright\arraybackslash}p{#1}}
\captionsetup{skip=7pt}
\setlength{\textfloatsep}{17pt plus 3pt minus 3pt}
\renewcommand{\bibfont}{\small}
\setlength{\bibsep}{4pt plus 1pt minus 1pt}
\setlist{itemsep=3pt,topsep=5pt}
\numberwithin{equation}{section}
\newtheorem{theorem}{Theorem}[section]
\newaliascnt{lemma}{theorem}
\newtheorem{lemma}[lemma]{Lemma}
\aliascntresetthe{lemma}
\newaliascnt{proposition}{theorem}
\newtheorem{proposition}[proposition]{Proposition}
\aliascntresetthe{proposition}
\newaliascnt{corollary}{theorem}
\newtheorem{corollary}[corollary]{Corollary}
\aliascntresetthe{corollary}
\theoremstyle{definition}
\newaliascnt{definition}{theorem}
\newtheorem{definition}[definition]{Definition}
\aliascntresetthe{definition}
\newaliascnt{algorithm}{theorem}
\newtheorem{algorithm}[algorithm]{Algorithm}
\aliascntresetthe{algorithm}
\theoremstyle{remark}
\newaliascnt{remark}{theorem}
\newtheorem{remark}[remark]{Remark}
\aliascntresetthe{remark}
\crefname{algorithm}{Algorithm}{Algorithms}
\newcommand{\R}{\mathbb R}
\newcommand{\N}{\mathbb N}
\newcommand{\E}{\mathbb E}
\newcommand{\PP}{\mathbb P}
\newcommand{\one}{\mathbf 1}
\newcommand{\cI}{\mathcal I}
\newcommand{\cD}{\mathcal D}
\newcommand{\cW}{\mathcal W}
\newcommand{\cH}{\mathcal H}
\newcommand{\cL}{\mathcal L}
\newcommand{\cB}{\mathcal B}
\newcommand{\cT}{\mathsf T}
\newcommand{\Tcoord}{T_{\mathrm{coord}}^\star}
\newcommand{\Tupd}{T_{\mathrm{upd}}^\star}
\newcommand{\Tbox}{T_{\Box}}
\newcommand{\Csel}{C_{\mathrm{sel}}}
\newcommand{\Cen}{C_{\mathrm{en}}}
\newcommand{\eps}{\epsilon}
\DeclareMathOperator{\vol}{vol}
\DeclareMathOperator{\aff}{aff}
\DeclareMathOperator{\argmin}{arg\,min}
\DeclareMathOperator{\argmax}{arg\,max}
\DeclareMathOperator{\supp}{supp}
\DeclareMathOperator{\diam}{diam}
\DeclareMathOperator{\sign}{sign}
\DeclareMathOperator{\rank}{rank}
\DeclareMathOperator{\ind}{ind}
\DeclareMathOperator{\conv}{conv}
\DeclarePairedDelimiter{\norm}{\lVert}{\rVert}
\DeclarePairedDelimiter{\abs}{\lvert}{\rvert}
\newcommand{\ip}[2]{\langle #1,#2\rangle}
\newcommand{\ball}{B_2^d}
\title{Near-Optimal Complexity of Convex Optimization\\with Exact Coordinate Subgradient Oracles}
\author{Colombo Team\\[3pt]
\small\href{mailto:colombo.2026@outlook.com}{\texttt{colombo.2026@outlook.com}}}
\date{}

\begin{document}
\maketitle
\vspace{-1.3em}
\begin{abstract}
We study convex Lipschitz optimization when each query returns one exact coordinate of a fixed bounded subgradient selection. The feasible set is a known compact product of intervals, and one optimizer is assumed to lie within a prescribed distance of the initial point. Algorithms may randomize, query arbitrary feasible points, and select coordinates adaptively; the number of scalar queries is bounded on every execution. Let $q=GR_B/\eps$ and $s=\min\{q^2,d\}$, where $G$ bounds the Euclidean subgradient norm and $\sqrt2R_B$ is the optimizer-distance bound. For all sufficiently large $d$ and $q$, we prove a worst-instance expected-error lower bound of $\Omega(ds/\log(es))$. The proof extends the posterior-mean energy argument for exact maximum observations to selected-coordinate feedback: a coordinate reveal imposes an additional affine slice with a dimension-free energy cost. A radial penalty permits unrestricted queries on the public box, and a direct expectation argument converts posterior energy into optimization error. Complementary upper bounds are $O(dq^2)$ for projected coordinate subgradient updates and $O(d^2\log(edq))$ for deterministic, value-free centroid localization. The latter also admits coordinatewise motion with at most twice as many scalar queries. These bounds identify the joint polynomial scale $d\min\{q^2,d\}$ up to logarithmic factors. The precise logarithmic dependence and sharp dimension-specific bounds are not determined by the present argument.
\end{abstract}
\noindent\textbf{Keywords:} convex optimization; coordinate subgradient oracle; randomized coordinate methods; oracle complexity; posterior-mean energy; cutting planes.

\medskip
\begin{samepage}
\noindent\textbf{AI Usage.}
Nearly the entire research pipeline for this paper was carried out by
Colombo, powered by GPT-5.6 Sol.
\par
\end{samepage}

\section{Introduction}\label{sec:intro}

Coordinate methods reduce the information used at each iteration by replacing a full subgradient with one of its coordinates. For a convex objective with a bounded subgradient selection $g$, uniform sampling gives the unbiased estimator
\[
H=d\,g_I(x)e_I,\qquad
\E[H\mid x]=g(x),\qquad
\E[\norm{H}_2^2\mid x]=d\norm{g(x)}_2^2.
\]
If $\norm{g(x)}_2\le G$ and an optimizer lies within distance $\sqrt2R_B$ of the initial point, the usual projected subgradient analysis yields expected error $O(GR_B\sqrt{d/T})$ after $T$ scalar queries. This explains the dimension factor in a familiar algorithm. Whether the same factor is unavoidable for arbitrary adaptive algorithms is a different question.

An unrestricted algorithm can query dense points, choose coordinates using all previous responses, return to the same point, or use global localization instead of a subgradient update. Exact observations make this distinction consequential: one real number can encode more than a fixed number of bits. For a nonsmooth objective, even the identity of the active subgradient at a dense query can depend on every hidden direction. Neither a bit-counting argument nor a lower bound for a prescribed update rule settles the complexity of the oracle itself.

We consider the following information interface. A hidden instance consists of a convex objective $f$ and a bounded Borel selection $g(x)\in\partial f(x)$, fixed before the algorithm's random seed is realized. A query $(x,i)$ returns only $g_i(x)$. All feasible points and coordinates may be chosen adaptively. The domain is a known compact product of intervals, and only the distance to one optimizer is controlled. In particular, the domain diameter can be much larger than the optimizer-distance scale. Throughout the paper, Lipschitz continuity refers to the objective; differentiability and gradient Lipschitz continuity are not assumed.

\subsection{Main result and accuracy regimes}

Write
\begin{equation}\label{eq:scales-intro}
q=\frac{GR_B}{\eps},\qquad s=\min\{q^2,d\}.
\end{equation}
Let $\Tcoord(d,G,R_B,\eps)$ denote the smallest pathwise scalar-query budget that guarantees worst-instance expected error at most $\eps$, uniformly over the public product geometries. For universal positive constants $c,C$, and all $d\ge d_0$ and $q\ge q_0$, our bounds are
\begin{equation}\label{eq:intro-frontier}
c\frac{ds}{\log(es)}
\le \Tcoord(d,G,R_B,\eps)
\le C\min\{dq^2,d^2\log(edq)\}.
\end{equation}
The lower bound already holds on one public box, with a canonical subgradient selection fixed as part of each instance. The upper bounds hold in every dimension and for every admissible product geometry. \Cref{tab:regimes} displays the two regimes without suppressing their logarithms.

\begin{table}[t]
\centering
\caption{Scalar-coordinate query bounds in this paper. In all rows, $d\ge d_0$ and $q\ge q_0$. The error criterion is worst-instance expected suboptimality and the budget is pathwise. Constants in $O(\cdot)$ and $\Omega(\cdot)$ are universal.}
\label{tab:regimes}
\small
\setlength{\tabcolsep}{5pt}
\renewcommand{\arraystretch}{1.35}
\begin{tabularx}{\linewidth}{@{}L{0.21\linewidth}L{0.25\linewidth}Y@{}}
\toprule
Accuracy regime & Lower bound & Upper bound and implementation \\
\midrule
$q_0\le q\le\sqrt d$
& $\displaystyle\Omega\!\left(\frac{dq^2}{\log(eq^2)}\right)$
& $O(dq^2)$, using projected coordinate subgradient updates. \\
$q\ge\sqrt d$
& $\displaystyle\Omega\!\left(\frac{d^2}{\log(ed)}\right)$
& $O(\min\{dq^2,d^2\log(edq)\})$, choosing the cheaper of coordinate updates and centroid localization. \\
\bottomrule
\end{tabularx}
\end{table}

At $q\asymp\sqrt d$, projected coordinate updates give an $O(d^2)$ upper bound. At higher accuracy, centroid localization can be preferable. The transition in the lower bound occurs at $\sqrt d$; the crossing of the two explicit upper bounds also involves a logarithm. In particular, the high-accuracy lower bound does not establish independence from $q$, because the upper bound still contains $\log q$.

Our contributions are the following.
\begin{enumerate}[leftmargin=*,label=\arabic*.]
\item \textbf{A lower bound for adaptive coordinate access.} We transfer the exact-maximum posterior-energy framework of Zhang et al.\ \cite{ZZQL2026} to one coordinate of the selected active row. The proof covers dense query points, adaptive coordinate choices, repeated queries, and arbitrary independent standard Borel random seeds.
\item \textbf{Upper bounds in the same resource model.} We give the standard projected coordinate subgradient bound under an aggregate Euclidean norm assumption, and a deterministic centroid policy that uses no function values. Both have explicit pathwise budgets. On a product domain, a coordinatewise routing lemma realizes the centroid policy while changing only the most recently queried coordinate between successive queries, with at most twice as many calls.
\item \textbf{A precise account of oracle consistency.} The proofs specify the fixed selection, the order of quantifiers, and the behavior on arbitrary response histories. An auxiliary appendix characterizes the convex Lipschitz completions of a finite coordinate transcript and explains why allowing the optimizer location to vary introduces a nonconvex feasibility condition.
\end{enumerate}

\subsection{Relation to prior work}\label{sec:related}

\paragraph{Coordinate algorithms.}
Nesterov\ \cite{Nesterov2012} established foundational complexity results for coordinate descent under coordinate smoothness. Nonsmooth block methods were developed by Dang and Lan\ \cite{DangLan2015} and Deng, Lan, and Rangarajan\ \cite{DengLanRangarajan2015}. Zhao et al.\ \cite{ZhaoEtAl2023} analyze randomized coordinate subgradient methods for nonsmooth composite objectives. Their bounded-subgradient recursion has the same $O(dq^2)$ fixed-horizon specialization as our coordinate upper bound; \cref{tab:comparison} specifies the difference in access models.

The $T^{-1/2}$ subgradient convergence mechanism is classical, but its assumptions and parameters matter when comparing dimension dependence. Bounds expressed through separate uniform block moments do not automatically reduce to a bound on the Euclidean norm of the full subgradient. Indeed,
\[
\sup_{x\in D}\sum_i|g_i(x)|^2\le G^2
\quad\not\Rightarrow\quad
\sum_i\sup_{x\in D}|g_i(x)|^2\le G^2.
\]
Our elementary proof uses the pointwise aggregate bound directly. Smooth-coordinate complexity results concern a different function class from the nonsmooth class studied here.

\paragraph{Exact scalar lower bounds.}
The closest lower-bound result is Zhang et al.\ \cite{ZZQL2026}, who establish the polynomial scale $d\min\{d,\eps^{-2}\}$ up to logarithms for exact function-value queries on a Euclidean ball at fixed Lipschitz and radius scales. Their central analytic input is a one-step bound on the posterior-mean displacement of a sum of independent strongly log-concave rows after a maximum value and its active label are revealed. We use that theorem in its stated affine-support form, including the one-row case.

Our coordinate oracle returns a component of the active row, whereas their oracle returns the maximum value. Neither observation can be identified with the other in the proof. Here the posterior analysis must account for an additional exact coordinate equality; queries are permitted throughout a product box whose diameter is not controlled by the optimizer radius; and the main guarantee is stated directly in expected error. These are the extensions established below. The random support-function construction and the underlying posterior-energy method are credited to Zhang et al.; we do not claim them as independent contributions.

\paragraph{Centroid localization and convex interpolation.}
The high-accuracy upper bound uses classical centroid geometry. The Minkowski--Radon asymmetry and Gr\"unbaum halfspace inequalities are used in the forms stated by Alonso-Guti\'errez et al.\ \cite{AGMSMGYN2024}; the shifted-cap argument follows Nasz\'odi\ \cite[Lemma~2.1]{Naszodi2017}. We include the geometric reduction, a support-gap stopping certificate, and the proof of a fixed query budget under the optimizer-radius promise. Exact centroid and support computations are free internal operations in the oracle model. The resulting query bound makes no assertion about arithmetic runtime. The finite-transcript appendix specializes classical max-affine interpolation and cyclic monotonicity\ \cite{Rockafellar1966} to partially observed coordinates of a fixed subgradient selection.

\begin{table}[t]
\centering
\caption{Comparison of related guarantees. The rows distinguish oracle access, geometry, and error criterion; they are not claims that the models are equivalent. In the exact-value row, Lipschitz and radius scales are normalized constants and $s=\min\{d,q^2\}$ with $q$ proportional to $\eps^{-1}$. In the other rows, $q=GR_B/\eps$ and $s=\min\{d,q^2\}$.}
\label{tab:comparison}
\footnotesize
\setlength{\tabcolsep}{4pt}
\renewcommand{\arraystretch}{1.3}
\begin{tabularx}{\linewidth}{@{}L{0.19\linewidth}L{0.31\linewidth}L{0.23\linewidth}Y@{}}
\toprule
Work & Setting and information & Bound & Error criterion \\
\midrule
Zhao et al.\ \cite{ZhaoEtAl2023}; bounded-subgradient specialization
& Unconstrained convex composite objective; a uniformly sampled coordinate subgradient is formed from composite structure.
& $O(dq^2)$ coordinate updates
& Expected suboptimality. \\
\addlinespace[4pt]
Zhang et al.\ \cite{ZZQL2026}
& Exact scalar function values on a Euclidean ball; arbitrary adaptive randomized algorithms.
& $\displaystyle\Omega\!\left(\frac{ds}{\log(es)}\right)$
& Success probability at least $1/2$. \\
\addlinespace[4pt]
This paper: lower bound
& Exact coordinate of a fixed bounded subgradient selection; one public product box.
& $\displaystyle\Omega\!\left(\frac{ds}{\log(es)}\right)$
& Worst-instance expected suboptimality. \\
\addlinespace[4pt]
This paper: coordinate updates
& Same scalar oracle; any admissible product domain and a nearby optimizer.
& $\lceil2dq^2\rceil$ scalar calls
& Expected suboptimality. \\
\addlinespace[4pt]
This paper: centroid localization
& Same scalar oracle and geometry; exact internal centroid and support computations.
& $O(d^2\log(edq))$ scalar calls
& Deterministic suboptimality. \\
\bottomrule
\end{tabularx}

\vspace{5pt}
\parbox{\linewidth}{\footnotesize
The Zhao et al.\ rate is the fixed-horizon specialization of \cite[Lemma~3.1(b)]{ZhaoEtAl2023} to scalar blocks, $\norm{g(x)}_2\le G$, and $\norm{x^0-x^\star}_2\le\sqrt2R_B$. It counts coordinate updates and is not a minimax result for the present scalar black-box model. The two lower-bound rows require sufficiently large $d$ and $q$; the upper bounds proved here hold for every $d\ge1$ and $q>1/\sqrt2$. Coordinatewise routing implements the last row with at most twice the scalar budget.}
\end{table}

\subsection{Proof overview}

The lower-bound family is a maximum of $K$ independent random linear functions plus a radial penalty. The penalty places every optimizer in the unit ball while keeping the entire public box available for queries. A fixed least-index rule selects one active row. Subtracting the known radial contribution from a scalar oracle response reveals exactly one coordinate of that row.

For analysis, we augment each response by the active-row label and the maximum score. Conditional on the augmented transcript, the hidden rows remain independent truncated Gaussian laws on affine supports. A label-and-score observation adds a winner equality and separate loser halfspaces. The observed coordinate adds one further equality to the winner row. The selection inequality of Zhang et al.\ \cite[Theorem~14]{ZZQL2026} bounds the first posterior-mean increment by $O(\log(eK))$ and its one-row case bounds the second by $O(1)$. Orthogonality of conditional-expectation increments then gives
\[
\E\left\|\E\left[\sum_{j=1}^K a_j\,\middle|\,\cW_T\right]\right\|_2^2
\lesssim \frac{T\log(eK)}d.
\]
The optimum is negative on the scale $K^{-1/2}$. On the other hand, a maximum is at least the average of its linear pieces, so a successful output must correlate with the negative aggregate direction. Conditional expectation and Cauchy--Schwarz convert the energy bound directly into an expected-error lower bound. The choice $K\asymp\min\{q^2,d\}$ gives both accuracy regimes.

\paragraph{Organization.}
\Cref{sec:model} defines the oracle model and states the main results. \Cref{sec:coordinate-upper,sec:cutting-plane} prove the upper bounds, including coordinatewise routing. \Cref{sec:hard,sec:posterior,sec:lower} develop the hard family, the adaptive posterior, and the lower bound. \Cref{sec:scope} discusses the scope of the results. The appendices provide measurable-kernel details, an optional probability-of-failure version, and finite-transcript interpolation.

\section{Oracle model and main theorems}\label{sec:model}

\subsection{Public geometry and hidden instance}

Write $[d]=\{1,\ldots,d\}$ and $\N_0=\{0,1,2,\ldots\}$. The notation $B_2^d$ denotes the closed Euclidean unit ball. For a convex set $C$, $\ind_C$ is its convex indicator, equal to zero on $C$ and $+\infty$ outside; $\one_C$ denotes its ordinary zero--one indicator. Let $\cD_d$ be the collection of known pairs $(D,x_1)$, where
\begin{equation}\label{eq:box-domain}
D=\prod_{i=1}^d[\ell_i,u_i],\qquad
-\infty<\ell_i\le u_i<\infty,\qquad x_1\in D.
\end{equation}
Singleton factors are allowed. There is no uniform bound on $\diam(D)/R_B$.

\begin{definition}[Coherent coordinate instance]\label{def:instance}
For $G,R_B>0$, the class $\cI(D,x_1;G,R_B)$ consists of pairs $(f,g)$ with the following properties. The function $f:D\to\R$ is finite, convex, and $G$-Lipschitz in the Euclidean norm. The map $g:D\to\R^d$ is one fixed Borel map satisfying, for all $x,z\in D$,
\begin{equation}\label{eq:coherence}
f(z)\ge f(x)+\ip{g(x)}{z-x},\qquad \norm{g(x)}_2\le G.
\end{equation}
There is at least one $x^\star\in\argmin_D f$ such that
\begin{equation}\label{eq:radius}
\frac12\norm{x^\star-x_1}_2^2\le R_B^2.
\end{equation}
The exact coordinate oracle returns $g_i(x)$ on query $(x,i)\in D\times[d]$.
\end{definition}

The term \emph{coherent} means that all coordinates at a given point come from the same vector, and that this selection is fixed before interaction. The selection is part of the hidden instance rather than a choice made retrospectively in response to the algorithm's randomness. At boundary points, \eqref{eq:coherence} is understood relative to $D$. All our arguments use this displayed inequality directly.

\subsection{Policies, randomness, and scalar accounting}

For a fixed public tuple $(d,G,R_B,\eps,D,x_1)$, a randomized policy uses a standard Borel seed space with an instance-independent probability law. Its seed is denoted by $U$. A length-$t$ history has the form
\[
H_t=((X_s,I_s,Y_s))_{s=1}^t\in
\cH_t:=(D\times[d]\times[-G,G])^t,
\qquad Y_s=g_{I_s}(X_s).
\]
The query $(X_{t+1},I_{t+1})$ and terminal output $\widehat x$ are jointly Borel functions of the seed and available history. These maps are defined and feasible on every possible history. The cap $T$ is pathwise: no seed realization or response path uses more than $T$ scalar responses. Early stopping can be represented by legal ignored queries up to the common cap.

The policy is chosen after the public geometry is known. The adversary then fixes the complete deterministic pair $(f,g)$, and the seed is realized last. The policy receives no function values, active labels, maximum scores, full vectors, coordinate minimizers, or off-domain responses. Arbitrary exact Borel internal computations are allowed at no oracle cost.

\begin{definition}[Worst-geometry coordinate complexity]\label{def:complexity}
With the convention $\inf\varnothing=+\infty$, define
\begin{equation}\label{eq:minimax}
\begin{split}
\Tcoord(d,G,R_B,\eps)=\inf\bigl\{T\in\N_0:\;&
\forall(D,x_1)\in\cD_d\ \exists\mathsf A_{D,x_1}\text{ with cap }T,\\
&\sup_{(f,g)\in\cI(D,x_1;G,R_B)}
\E_U[f(\widehat x)-\min_Df]\le\eps\bigr\}.
\end{split}
\end{equation}
\end{definition}
Only the objective gap is averaged. Feasibility and the query cap are not average-case requirements.

On the normalized box $D_d=[-1,1]^d$, with $x_1=0$ and $G=R_B=1$, write
\begin{equation}\label{eq:tbox}
\Tbox(d,q)=\inf\left\{T\in\N_0:
\exists\mathsf A\ \sup_{(f,g)\in\cI(D_d,0;1,1)}
\E_U[f(\widehat x)-\min f]\le\frac1q\right\}.
\end{equation}

\subsection{Main statements}

\begin{theorem}[Coordinate-oracle frontier]\label{thm:main}
There are universal constants $c>0$, $q_0\ge1$, and $d_0\in\N$ such that, for every $d\ge d_0$, $G,R_B>0$, and $q=GR_B/\eps\ge q_0$, with $s=\min\{q^2,d\}$,
\begin{equation}\label{eq:main-lower}
\Tcoord(d,G,R_B,\eps)\ge c\frac{ds}{\log(es)}.
\end{equation}
The normalized lower bound $\Tbox(d,q)\ge cds/\log(es)$ already holds on the single public box $[-1,1]^d$.

For every $d\ge1$, $G,R_B>0$, and $q>1/\sqrt2$,
\begin{equation}\label{eq:main-upper}
\Tcoord(d,G,R_B,\eps)\le
\min\left\{\lceil2dq^2\rceil,
\ d\bigl(\lceil2ed\log(4\sqrt2\,d^2q)\rceil+1\bigr)\right\}.
\end{equation}
The first bound is attained by a randomized one-coordinate update method with an averaged output. The second is attained deterministically without function values. In particular, the second budget is $O(d^2\log(edq))$ for $q\ge1$.
\end{theorem}

\begin{definition}[Coordinate-update policies]\label{def:update}
A coordinate-update policy is a coordinate-oracle policy whose first query point is $x_1$ and whose actual successive query points satisfy
\[
X_{t+1}-X_t\in\operatorname{span}\{e_{I_t}\}.
\]
The output may be any feasible Borel function of the transcript, including an average of queried iterates. Let $\Tupd$ be defined as in \eqref{eq:minimax} with this restricted policy class.
\end{definition}

\begin{corollary}[Complexity under coordinatewise motion]\label{cor:update}
For every $d\ge1$ and $q>1/\sqrt2$,
\begin{equation}\label{eq:update-upper}
\Tcoord\le\Tupd\le
\min\left\{\lceil2dq^2\rceil,\,
2d\bigl(\lceil2ed\log(4\sqrt2\,d^2q)\rceil+1\bigr)\right\},
\end{equation}
where the suppressed arguments are $(d,G,R_B,\eps)$. Consequently, for $d\ge d_0$ and $q\ge q_0$,
\begin{equation}\label{eq:update-frontier}
c\frac{d\min\{q^2,d\}}{\log(e\min\{q^2,d\})}
\le \Tupd(d,G,R_B,\eps)
\le C\min\{dq^2,d^2\log(edq)\}
\end{equation}
for a universal $C$. In the regime $q_0\le q\le\sqrt d$, the upper bound is attained by the projected coordinate subgradient method. In the high-accuracy regime, the second implementation uses centroid computations and coordinatewise routing.
\end{corollary}
\begin{proof}
The lower bound follows because coordinate-update policies form a subset of all coordinate-oracle policies. \Cref{thm:rcs} gives the first upper bound; \cref{cor:routed-centroid} gives the second. The simplified bound follows from \cref{cor:cp-log}.
\end{proof}

\begin{remark}[A motion restriction is not an update formula]
\Cref{def:update} constrains successive query points, but permits arbitrary Borel internal computations and arbitrary feasible outputs. A prescribed subgradient recursion is a narrower class. The high-accuracy guarantee in \cref{cor:update} concerns the stated motion restriction; it is not a convergence theorem for every coordinate-update rule or for its last iterate.
\end{remark}

\begin{proposition}[Zero-query regimes]\label{prop:zero}
If $\eps\ge\sqrt2GR_B$, outputting $x_1$ without queries is sufficient. The same is true, with zero objective error, if $G=0$ or $R_B=0$ in the natural extension of the instance class.
\end{proposition}
\begin{proof}
By Lipschitz continuity and \eqref{eq:radius}, $f(x_1)-f(x^\star)\le\sqrt2GR_B$. If $G=0$, $f$ is constant. If $R_B=0$, the radius promise supplies $x^\star=x_1$.
\end{proof}

\begin{remark}[What the dimension parameter measures]
The ambient dimension $d$ is the number of scalar subgradient coordinates. Our upper bounds remain valid when the effective affine dimension of $D$ is smaller. The lower bound uses a full-dimensional box. No lower bound for every individual product geometry is claimed; a singleton domain, for example, is trivial.
\end{remark}

\section{Projected coordinate subgradient updates}\label{sec:coordinate-upper}

The following proof is included to make the resource model and update restriction explicit. In particular, we project onto the product $D$, not onto its intersection with a Euclidean ball. The latter projection could alter coordinates other than the sampled one.

\begin{algorithm}[Uniform projected coordinate subgradient method]\label{alg:rcs}
Fix an integer $N\ge1$ and set
\begin{equation}\label{eq:rcs-eta}
\eta=\frac{\sqrt2R_B}{G\sqrt{dN}}.
\end{equation}
Draw independent uniform coordinates $I_1,\ldots,I_N\in[d]$, independently of the instance. Start at $X_1=x_1$. At time $t$, query $(X_t,I_t)$, receive $g_{I_t}(X_t)$, and define
\begin{equation}\label{eq:rcs-update}
H_t=d\,g_{I_t}(X_t)e_{I_t},\qquad
X_{t+1}=\Pi_D(X_t-\eta H_t).
\end{equation}
Return $\widehat x=N^{-1}\sum_{t=1}^N X_t$. The auxiliary point $X_{N+1}$ is not queried.
\end{algorithm}

For a box, $\Pi_D$ is coordinatewise clipping. Since $X_t\in D$, every coordinate except $I_t$ remains unchanged in \eqref{eq:rcs-update}.

\begin{lemma}[Projection inequality]\label{lem:projection}
For a nonempty closed convex set $C\subset\R^d$, its Euclidean projection satisfies
\begin{equation}\label{eq:projection}
\norm{\Pi_C(y)-z}_2^2\le\norm{y-z}_2^2
\qquad(y\in\R^d,\ z\in C).
\end{equation}
For the compact sets used here, projection is single-valued, continuous, and $1$-Lipschitz.
\end{lemma}
\begin{proof}
Existence for a compact set follows by minimization of squared distance, and uniqueness follows from strict convexity. Let $p=\Pi_C(y)$. Comparing $p$ with $p+t(z-p)$ and letting $t\downarrow0$ yields $\ip{y-p}{z-p}\le0$. Expanding $\norm{y-z}_2^2$ proves \eqref{eq:projection}. Applying the variational inequality to $p=\Pi_C(y)$ and $p'=\Pi_C(y')$ and adding gives
\[
\norm{p-p'}_2^2\le\ip{p-p'}{y-y'}\le\norm{p-p'}_2\norm{y-y'}_2.
\]
Thus projection is $1$-Lipschitz. Existence for general nonempty closed convex $C$ follows by restricting the distance minimization to a sufficiently large closed ball.
\end{proof}

\begin{theorem}[Coordinate risk bound]\label{thm:rcs}
For every admissible geometry and hidden pair, \cref{alg:rcs} is a legal $N$-query coordinate-update policy and satisfies
\begin{equation}\label{eq:rcs-risk}
\E_U[f(\widehat x)-\min_D f]
\le\sqrt2GR_B\sqrt{\frac dN}.
\end{equation}
Consequently, $N=\lceil2dq^2\rceil$ attains error at most $\eps$.
\end{theorem}
\begin{proof}
Choose the optimizer in \eqref{eq:radius}. With $\mathcal F_{t-1}=\sigma(I_1,\ldots,I_{t-1})$, the iterate $X_t$ is $\mathcal F_{t-1}$-measurable, and
\begin{align}
\E[H_t\mid\mathcal F_{t-1}]&=g(X_t),\label{eq:unbiased}\\
\E[\norm{H_t}_2^2\mid\mathcal F_{t-1}]&=d\norm{g(X_t)}_2^2\le dG^2.\label{eq:secondmoment}
\end{align}
By \cref{lem:projection},
\[
\norm{X_{t+1}-x^\star}_2^2
\le \norm{X_t-x^\star}_2^2
-2\eta\ip{H_t}{X_t-x^\star}+\eta^2\norm{H_t}_2^2.
\]
Taking conditional expectations and using \eqref{eq:coherence}, followed by summation, gives
\begin{equation}\label{eq:rcs-telescope}
\frac1N\sum_{t=1}^N\E[f(X_t)-f(x^\star)]
\le \frac{\norm{x_1-x^\star}_2^2}{2\eta N}
+\frac{\eta dG^2}{2}
\le\frac{R_B^2}{\eta N}+\frac{\eta dG^2}{2}.
\end{equation}
Convexity bounds the expected error of the average by the left side. Substituting \eqref{eq:rcs-eta} proves \eqref{eq:rcs-risk}.

For completeness, the seed space is the finite standard Borel space $[d]^N$. At each possible history, use the last displayed point, coordinate, and scalar response in the continuous clipping formula. The next point is feasible, even on a history inconsistent with an actual instance. The terminal average is feasible by convexity and is Borel. All query maps are Borel, the number of scalar responses is exactly $N$, and the updates change only the coordinate most recently queried. No bound on the distance of later iterates from $x_1$ is needed in \eqref{eq:rcs-telescope}.
\end{proof}

\begin{remark}[A queried-point output]
Instead of averaging, the policy may independently choose a uniform index $J\in[N]$ and return $X_J$. Equation~\eqref{eq:rcs-telescope} gives the same expected-error bound. Thus the upper bound can also be realized with an output that is one of the queried iterates. Neither formulation is a bound on the final iterate $X_N$.
\end{remark}

\section{Value-free centroid localization}\label{sec:cutting-plane}

We next prove the second part of \cref{thm:main}. The algorithm reconstructs one coherent subgradient at a localization center using $d$ scalar calls. It uses a computable support-gap certificate, rather than function values, to decide when to stop. A fixed shallow offset preserves a positive-volume set of near-optimal points and gives a deterministic cap.

\subsection{Localization using only the optimizer-radius promise}

Let
\begin{equation}\label{eq:localization}
r=\sqrt2R_B,\qquad
J=\{i:\ell_i<u_i\},\qquad n=|J|,
\end{equation}
\[
E=\{x\in\R^d:x_i=(x_1)_i\text{ for }i\notin J\},
\qquad K_0=D\cap\overline B_2(x_1,r).
\]
All centroids, volumes, and relative interiors below are taken in $E$, identified with $\R^n$ through its active coordinates.

\begin{lemma}[Initial localization]\label{lem:initial}
If $n=0$, $D=\{x_1\}$. If $n\ge1$, $K_0$ is a compact convex body with positive relative $n$-volume. It contains an optimizer $x^\star$ satisfying \eqref{eq:radius}, and
\begin{equation}\label{eq:local-loss}
0\le f(z)-f(x^\star)\le2rG\qquad(z\in K_0).
\end{equation}
\end{lemma}
\begin{proof}
Each active interval contains a one-sided interval of positive length starting at $(x_1)_i$. A sufficiently small product of these intervals lies in $D\cap\overline B_2(x_1,r)$, proving positive relative volume. Compactness and convexity are immediate. The promised optimizer lies in $K_0$. Both it and $z$ are within $r$ of $x_1$, so Lipschitz continuity yields \eqref{eq:local-loss}.
\end{proof}

\subsection{Centroid geometry and shallow cuts}

We isolate the external geometric facts. If a full-dimensional compact convex body $L\subset\R^n$ has centroid zero, then
\begin{equation}\label{eq:centroid-asymmetry}
L\subseteq-nL,
\end{equation}
and every halfspace through zero contains at least $\vol_n(L)/e$ of its volume. These are the Minkowski--Radon and Gr\"unbaum inequalities, respectively; see \cite[Theorem~A and equation~(1.1)]{AGMSMGYN2024}. They apply after translation and identification with the relative affine hull. In dimension one they follow directly from the midpoint of an interval.

\begin{lemma}[Shifted cap and contraction]\label{lem:cap}
Let $K$ be a convex body of relative dimension $n\ge1$ with centroid $c$. For a linear functional $v$, set
\begin{equation}\label{eq:supports}
A=\max_{y\in K}\ip{v}{c-y},\qquad
B=\max_{y\in K}\ip{v}{y-c}.
\end{equation}
If $A>\eps$ and $\delta=\eps/(2n^2)$, then
\begin{equation}\label{eq:cap-contraction}
\vol_n\bigl(K\cap\{y:\ip{v}{y-c}\le\delta\}\bigr)
\le\rho\vol_n(K),\qquad \rho=1-\frac1{2e}.
\end{equation}
\end{lemma}
\begin{proof}
Translate $c$ to zero and write $L=K-c$. The active restriction of $v$ is nonzero because $A>0$. From \eqref{eq:centroid-asymmetry},
\[
B=h_L(v)\ge\frac{h_L(-v)}n=\frac An>\frac\eps n.
\]
Thus $\vartheta=\delta/B$ belongs to $(0,1/(2n))$. Choose $p\in L$ with $\ip{v}{p}=B$, and let $L_+=L\cap\{z:\ip{v}{z}\ge0\}$. By convexity,
\[
\vartheta p+(1-\vartheta)L_+
\subseteq L\cap\{z:\ip{v}{z}\ge\delta\}.
\]
The left side has volume $(1-\vartheta)^n\vol_n(L_+)$. Gr\"unbaum's inequality and Bernoulli's inequality therefore imply
\[
\vol_n\bigl(L\cap\{z:\ip{v}{z}\ge\delta\}\bigr)
\ge\frac{(1-\vartheta)^n}{e}\vol_n(L)
\ge\frac{1-n\vartheta}{e}\vol_n(L)
>\frac1{2e}\vol_n(L).
\]
The boundary hyperplane has zero relative volume. Removing this cap proves \eqref{eq:cap-contraction}. The homothetic inclusion is precisely the shifted-cap argument in \cite[Lemma~2.1]{Naszodi2017}.
\end{proof}

\begin{lemma}[Value-free certificate and cut validity]\label{lem:certificate}
Suppose $K\subseteq K_0$ contains an optimizer $x^\star$, let $c$ be its relative centroid, and reconstruct $v=g(c)$ using the $d$ scalar calls $(c,1),\ldots,(c,d)$. Define
\begin{equation}\label{eq:certificate}
\mathcal A(K,c,v)=\max_{y\in K}\ip{v}{c-y}.
\end{equation}
If $\mathcal A(K,c,v)\le\eps$, then $f(c)-f(x^\star)\le\eps$. For every $\delta>0$, the halfspace
\begin{equation}\label{eq:shallow-cut}
H(c,v,\delta)=\{y\in E:\ip{v}{y-c}\le\delta\}
\end{equation}
contains every point $y\in K_0$ with $f(y)-f(x^\star)\le\delta$.
\end{lemma}
\begin{proof}
The subgradient inequality gives
\[
f(c)-f(x^\star)\le\ip{g(c)}{c-x^\star}\le\mathcal A(K,c,v).
\]
For a point with error at most $\delta$, the same inequality at $c$ gives
\[
\ip{v}{y-c}\le f(y)-f(c)\le f(y)-f(x^\star)\le\delta.
\]
These arguments use no observed or computed function value. Coherence ensures that the $d$ returned coordinates assemble to the single valid vector $g(c)$.
\end{proof}

\begin{lemma}[A fixed invariant core]\label{lem:core}
Assume $q>1/\sqrt2$ and $n\ge1$. Set
\begin{equation}\label{eq:core-parameters}
\delta=\frac\eps{2n^2},\qquad
\lambda=\frac\delta{2rG}=\frac1{4\sqrt2\,n^2q}.
\end{equation}
Then $0<\lambda<1$, and
\begin{equation}\label{eq:invariant-core}
Q=(1-\lambda)x^\star+\lambda K_0
\end{equation}
is contained in the $\delta$-sublevel set of $f-f(x^\star)$ within $K_0$. Moreover,
\begin{equation}\label{eq:core-volume}
\vol_n(Q)=\lambda^n\vol_n(K_0)>0.
\end{equation}
Every cut \eqref{eq:shallow-cut} generated at a feasible center preserves $Q$.
\end{lemma}
\begin{proof}
The condition on $q$ gives $\lambda<1/(4n^2)<1$. For $z\in K_0$, convexity and \eqref{eq:local-loss} yield
\[
f((1-\lambda)x^\star+\lambda z)-f(x^\star)
\le\lambda(f(z)-f(x^\star))\le2rG\lambda=\delta.
\]
The homothety has volume factor $\lambda^n$. Cut preservation follows from \cref{lem:certificate}.
\end{proof}

\subsection{The complete policy}

\begin{algorithm}[Fixed-budget shallow-centroid localization]\label{alg:centroid}
Assume $q>1/\sqrt2$ and define
\begin{equation}\label{eq:cp-budget}
M=\left\lceil2ed\log(4\sqrt2\,d^2q)\right\rceil+1,
\qquad T_{\mathrm{cp}}=dM.
\end{equation}
If $n=0$, store $x_1$. Otherwise start with body $K_0$ and a live state. At live block $k=1,\ldots,M$:
\begin{enumerate}[label=\arabic*.]
\item Compute the relative centroid $c_k$ of $K_{k-1}$.
\item Query $(c_k,1),\ldots,(c_k,d)$, keeping the point fixed throughout the block, and assemble the answers into $v_k$.
\item Compute $A_k=\mathcal A(K_{k-1},c_k,v_k)$. If $A_k\le\eps$, store $c_k$ and end the live state.
\item Otherwise form $\widetilde K_k=K_{k-1}\cap H(c_k,v_k,\delta)$. If $k<M$ and $\widetilde K_k$ has positive relative $n$-volume, set $K_k=\widetilde K_k$ and continue. On every other response history, enter an emergency state and store $x_1$.
\end{enumerate}
After a point is stored, fill all remaining scalar-response slots with the ignored legal query $(x_1,1)$. Return the stored point at cap $T_{\mathrm{cp}}$.
\end{algorithm}

\begin{theorem}[Termination and accuracy]\label{thm:centroid}
On every admissible instance, \cref{alg:centroid} reaches its accuracy test by block $M$, never reaches an emergency state, and returns a feasible point with
\begin{equation}\label{eq:cp-accuracy}
f(\widehat x)-\min_D f\le\eps
\end{equation}
deterministically. The scalar cap is \eqref{eq:cp-budget}.
\end{theorem}
\begin{proof}
The case $n=0$ is immediate. In the other case, \cref{lem:core} shows that every actual body and candidate cut contains $Q$. Thus its relative volume remains positive, and an optimizer remains in the body. Every nonterminal block contracts volume by \cref{lem:cap}. Consequently, after $k$ nonterminal blocks,
\[
\vol_n(K_k)\le\rho^k\vol_n(K_0),\qquad
\rho\le\exp\left(-\frac1{2e}\right).
\]
The function $u\mapsto u\log(4\sqrt2\,u^2q)$ is increasing for $u\ge1$ in the stated range of $q$. Hence
\[
M>2ed\log(4\sqrt2\,d^2q)
\ge2en\log(4\sqrt2\,n^2q)
=2en\log(1/\lambda).
\]
If the first $M$ blocks were all nonterminal, the last candidate body would contain $Q$, but would satisfy
\[
\vol_n(\widetilde K_M)
\le e^{-M/(2e)}\vol_n(K_0)
<\lambda^n\vol_n(K_0)=\vol_n(Q),
\]
a contradiction. Thus the accuracy test succeeds and \cref{lem:certificate} proves \eqref{eq:cp-accuracy}. Each full block costs exactly $d$ scalar calls; padding makes the cap pathwise.
\end{proof}

\begin{lemma}[Borel implementation and off-instance histories]\label{lem:cp-borel}
The construction in \cref{alg:centroid} defines feasible Borel query and output maps on every possible history, including histories inconsistent with any admissible selection.
\end{lemma}
\begin{proof}
Work in the known active-coordinate identification $E\simeq\R^n$, and fix a bounded box containing $K_0$. At every block boundary the current body is specified by membership in $K_0$ and finitely many linear inequalities whose coefficients are Borel functions of previous responses. Its indicator is jointly Borel in the history and integration variable. The volume and first moments are Borel parameter integrals. On the positive-volume branch their ratio is the centroid, hence a Borel map.

To show that $A_k$ is Borel without making a measurable selection of a support point, enumerate $\mathbb Q^n$. Rational points in a full-dimensional convex body are dense in it: approximate interior points by rational points and then approximate boundary points by segments from one fixed interior point. Therefore the maximum in \eqref{eq:certificate} equals the supremum of $\ip{v_k}{c_k-z}$ over the rational points $z$ lying in the body. Assign the value $-\infty$ when a rational point does not belong. This is a countable supremum of Borel functions.

The centroid of a compact convex body lies in the body; otherwise strict separation would contradict its defining average. Thus all live queries are feasible. The positive-volume test, certificate comparison, finite state transitions, and padding rules are Borel. Emergency outputs and padding queries use the known feasible point $x_1$. Within a block, the center is retained while the remaining coordinates are queried; intermediate scalar responses need not change the center. This gives the required Borel maps at every scalar-response time. Emergency branches ensure total definition, even though \cref{thm:centroid} shows that they are never reached on an admissible instance.
\end{proof}

\begin{corollary}[Simplified query budget]\label{cor:cp-log}
For $q\ge1$,
\begin{equation}\label{eq:cp-simple}
T_{\mathrm{cp}}
\le 2ed^2\log(4\sqrt2\,d^2q)+2d
\le(4e+2)d^2\log(edq).
\end{equation}
\end{corollary}
\begin{proof}
Use $\lceil a\rceil+1\le a+2$. Since $d,q\ge1$,
\[
\log(4\sqrt2\,d^2q)
\le2+2\log d+2\log q=2\log(edq).
\]
Also $d\le d^2\log(edq)$.
\end{proof}

\begin{remark}[Query complexity and computation]
The budget contains a single logarithm in $d$ and $q$. It follows from the deterministic number of centroid blocks, without assuming an efficient procedure for computing a centroid or a support maximum. The result concerns oracle calls only.
\end{remark}

\subsection{Implementing the policy with coordinatewise motion}

The product geometry permits a simple simulation that separates the motion of the query point from the internal computation used to choose it.

\begin{lemma}[Coordinatewise routing]\label{lem:routing}
Let $p,c\in D$, with $c$ determined by the available history. Starting at $p$, after $d$ ignored coordinate-oracle calls the next call may be made at $c$, while each successive query point differs from the preceding one only in the coordinate just queried. The route is feasible and Borel in $(p,c)$.
\end{lemma}
\begin{proof}
Set $p^0=p$ and, for $i=1,\ldots,d$, define
\[
p^i=(c_1,\ldots,c_i,p_{i+1},\ldots,p_d).
\]
Query $(p^{i-1},i)$, ignore its response, and use $p^i$ as the next query point. The displacement $p^i-p^{i-1}$ belongs to $\operatorname{span}\{e_i\}$. Since $D$ is a product and both endpoints lie in $D$, every intermediate point lies in $D$. Coordinate replacement is continuous, and $p^d=c$. The intermediate points need not belong to the localization body; this is allowed because queries throughout $D$ are legal.
\end{proof}

\begin{corollary}[Centroid localization with coordinatewise motion]\label{cor:routed-centroid}
There is a coordinate-update policy satisfying the deterministic accuracy guarantee \eqref{eq:cp-accuracy} with pathwise cap $2dM$, where $M$ is defined in \eqref{eq:cp-budget}.
\end{corollary}
\begin{proof}
If $n=0$, output $x_1$ and pad at $x_1$. Otherwise maintain the localization state of \cref{alg:centroid}, using only its designated center observations to update that state. Before each live block, compute the next center and apply \cref{lem:routing} from the current point, initially $x_1$. Then query all $d$ coordinates at that center. A block therefore uses $d$ routing calls and $d$ informative calls. The informative transcript is identical to that of \cref{alg:centroid}, so \cref{thm:centroid} proves accuracy within $M$ blocks.

At the boundary between blocks, the first routing query is made at the previous center, so the displacement after the preceding query is zero. Once the policy stores an output, all remaining calls are made at the current point and ignored. This preserves the motion restriction even when the stored output differs from the current point. On inconsistent histories use the same emergency branches as in \cref{alg:centroid}, with current-point padding. The localization maps are Borel by \cref{lem:cp-borel}; the routes are Borel by \cref{lem:routing}. Thus the policy is defined and feasible on every history, begins at $x_1$, and has the required cap.
\end{proof}

The simulation uses arbitrary internal localization computations. Its role is to show that a restriction on successive query-point motion, by itself, does not exclude the high-accuracy upper bound.

\section{A radially guarded hard distribution}\label{sec:hard}

We now work on the fixed normalized box $D_d=[-1,1]^d$, with $G=R_B=1$ and $x_1=0$. Let $K$ be an integer to be chosen later. Draw independent rows
\begin{equation}\label{eq:hard-law}
b_j\sim\cL\bigl(N(0,I_d)\mid\norm{N(0,I_d)}_2\le2\sqrt d\bigr),
\qquad a_j=\frac{b_j}{\sqrt d},\qquad j\in[K].
\end{equation}
This product law depends only on public parameters and is independent of the algorithmic seed. Let $A\in\R^{K\times d}$ have rows $a_j^\top$, and define
\begin{equation}\label{eq:support-function}
h_A(x)=\max_{j\in[K]}\ip{a_j}{x},\qquad
J_A(x)=\min\argmax_{j\in[K]}\ip{a_j}{x}.
\end{equation}
For the radial penalty and its selected subgradient, put
\begin{equation}\label{eq:radial-functions}
\rho(x)=(\norm{x}_2-1)_+,\qquad
r(x)=\begin{cases}
0,&\norm{x}_2\le1,\\
x/\norm{x}_2,&\norm{x}_2>1.
\end{cases}
\end{equation}
Our hidden pair is
\begin{equation}\label{eq:hard-pair}
F_A(x)=\frac{h_A(x)+4\rho(x)}6,\qquad
G_A(x)=\frac{a_{J_A(x)}+4r(x)}6.
\end{equation}
The ordering of the rows, least-index tie rule, and zero selection at the radial hinge are all fixed with the instance.

\subsection{Admissibility and radial projection}

\begin{lemma}[Admissibility and joint measurability]\label{lem:hard-borel}
For every ordered row array in the support of \eqref{eq:hard-law}, $F_A$ is finite, convex, and one-Lipschitz, and $G_A$ is a Borel subgradient selection with $\norm{G_A(x)}_2\le1$. The map $(A,x)\mapsto(F_A(x),G_A(x))$ is jointly Borel. For every admissible Borel policy, its finite transcript, terminal output, and terminal loss are jointly Borel in the finite row array and independent seed.
\end{lemma}
\begin{proof}
The row law has density proportional to
\[
e^{-\norm{b}_2^2/2}\one_{\{\norm{b}_2\le2\sqrt d\}}
\]
with a finite positive normalizing constant. It is a Borel probability law on a compact set. The function $h_A$ is a finite maximum of linear functions and is two-Lipschitz because $\norm{a_j}_2\le2$. The radial function $\rho$ is convex and one-Lipschitz. Thus $F_A$ is convex and one-Lipschitz.

The least-index selection cells satisfy
\begin{equation}\label{eq:tie-cells}
\{(A,x):J_A(x)=j\}
=\bigcap_{\ell<j}\{\ip{a_j-a_\ell}{x}>0\}
\cap\bigcap_{\ell>j}\{\ip{a_j-a_\ell}{x}\ge0\}.
\end{equation}
These are Borel sets, including at ties and duplicated rows. The selected active slope is a subgradient of $h_A$. The map $r$ in \eqref{eq:radial-functions} is a Borel subgradient of $\rho$; at $\norm{x}_2=1$, zero belongs to the radial subdifferential. Hence $G_A(x)\in\partial F_A(x)$, and
\[
\norm{G_A(x)}_2\le\frac{2+4}6=1.
\]
This proves the subgradient inequality on $D_d$.

The function $(A,x)\mapsto F_A(x)$ is continuous. Moreover, $A\mapsto\min_{D_d}F_A$ is continuous, because convergence of the finite row array gives uniform convergence on the compact box. Joint Borelness of the selection follows from \eqref{eq:tie-cells}. Induction through the policy's query maps makes every response and query jointly Borel in $(A,U)$; the terminal map and the continuous optimum-value map then make the terminal gap jointly Borel. No measurable structure on the abstract class of all functions is needed.
\end{proof}

\begin{lemma}[The radial guard]\label{lem:guard}
The continuous map
\begin{equation}\label{eq:radial-projection}
P(x)=\frac{x}{\max\{1,\norm{x}_2\}}
\end{equation}
maps $D_d$ into $\ball\subset D_d$ and satisfies $F_A(Px)\le F_A(x)$. Every minimizer of $F_A$ over $D_d$ lies in $\ball$, and
\begin{equation}\label{eq:guard-min}
\min_{D_d}F_A=\frac16\min_{\ball}h_A.
\end{equation}
In particular, the nearby-optimizer promise holds for every row array.
\end{lemma}
\begin{proof}
There is nothing to prove when $\norm{x}_2\le1$. Otherwise write $x=tu$ with $t>1$ and $\norm{u}_2=1$. Positive homogeneity gives
\[
F_A(tu)-F_A(u)=\frac{t-1}6\bigl(h_A(u)+4\bigr).
\]
Since every row has norm at most two, $h_A(u)\ge-2$, so this difference is strictly positive. A minimizer exists by compactness. It must lie in the unit ball, where the radial penalty vanishes, proving \eqref{eq:guard-min}. Its squared distance from zero divided by two is at most $1/2\le1$.
\end{proof}

\begin{remark}[Queries are not projected]
The projection in \cref{lem:guard} is used only to analyze the terminal loss. The policy may query any point of $[-1,1]^d$, including points outside the unit ball. Projecting all queries would improperly restrict the lower-bound model and is not done here.
\end{remark}

\subsection{An elementary spectral event}

The next lemma gives the spectral estimate used in \cite[Theorem~40]{ZZQL2026} in the range $K\le d/256$, with an elementary proof. This range is more than sufficient for our eventual choice $K=\lfloor a\min\{q^2,d\}\rfloor$ with a small universal $a$.

\begin{lemma}[Spectral event for the exact truncated law]\label{lem:spectral}
If $d\ge1024$ and $1\le K\le d/256$, the rows in \eqref{eq:hard-law} satisfy
\begin{equation}\label{eq:spectral}
\PP\left\{AA^\top\succeq\frac14 I_K\right\}\ge\frac78.
\end{equation}
\end{lemma}
\begin{proof}
Let $Z\in\R^{K\times d}$ have independent standard Gaussian entries, and set $H=ZZ^\top/d-I_K$. A Euclidean $1/4$-net $\mathcal N$ of the sphere in $\R^K$ can be chosen with $|\mathcal N|\le9^K$: take a maximal separated set and compare the disjoint balls of radius $1/8$ with a ball of radius $9/8$.

For any symmetric $H$, if $\mathcal N$ is such a net, then
\begin{equation}\label{eq:net-operator}
\norm{H}_{\mathrm{op}}\le2\max_{u\in\mathcal N}|u^\top Hu|.
\end{equation}
Indeed, choose a unit vector $x$ maximizing $|x^\top Hx|$ and $u\in\mathcal N$ within $1/4$ of $x$. The difference of the two quadratic forms is at most $2\norm{H}_{\mathrm{op}}\norm{x-u}_2\le\norm{H}_{\mathrm{op}}/2$.

For fixed $u$, $\norm{Z^\top u}_2^2$ has a chi-squared distribution with $d$ degrees of freedom. Its moment generating function is $(1-2t)^{-d/2}$ for $t<1/2$. Chernoff's bound yields, for $0<\alpha<1$,
\begin{equation}\label{eq:chisq}
\PP\left\{\left|\frac{\norm{Z^\top u}_2^2}{d}-1\right|>\alpha\right\}
\le2\exp\left(-\frac{d\alpha^2}{8}\right).
\end{equation}
For clarity, the upper-tail exponent is $d(\alpha-\log(1+\alpha))/2\ge d\alpha^2/8$, and the lower-tail exponent is $d(-\alpha-\log(1-\alpha))/2\ge d\alpha^2/4$.

Use $\alpha=3/8$, a union bound, and \eqref{eq:net-operator}. Since $K\le d/256$ and $\log9<5/2$,
\begin{align*}
\PP\{ZZ^\top/d\not\succeq I_K/4\}
&\le 2\exp\left(K\log9-\frac{9d}{512}\right)\\
&\le2e^{-d/128}.
\end{align*}
Let $\mathcal E$ be the event that every row of $Z$ has norm at most $2\sqrt d$. A chi-squared Chernoff bound at level $4d$ gives
\[
\PP(\mathcal E^c)\le K\exp\left(-\frac{3-\log4}{2}d\right)
\le d e^{-d/2}\le\frac12
\]
for the stated dimensions. Conditional on $\mathcal E$, the rows are independent with exactly the law \eqref{eq:hard-law}, before division by $\sqrt d$. Therefore
\[
\PP\{AA^\top\not\succeq I_K/4\}
\le\frac{2e^{-d/128}}{\PP(\mathcal E)}
\le4e^{-d/128}\le\frac18,
\]
where the last inequality holds for $d\ge1024$. This proves \eqref{eq:spectral}; the conditioning does not substitute an untruncated hard law.
\end{proof}

\begin{lemma}[Negative optimal value]\label{lem:negative-optimum}
Under the assumptions of \cref{lem:spectral}, if $h_A^\star=\min_{\ball}h_A$, then
\begin{equation}\label{eq:negative-optimum}
\PP\left\{h_A^\star\le-\frac1{4\sqrt K}\right\}\ge\frac78,
\qquad
\E h_A^\star\le-\frac7{32\sqrt K}.
\end{equation}
\end{lemma}
\begin{proof}
On the spectral event, set $\gamma=1/4$ and
\[
x_A=-\frac\gamma{\sqrt K}A^\top(AA^\top)^{-1}\one_K.
\]
Then $Ax_A=-(\gamma/\sqrt K)\one_K$ and
\[
\norm{x_A}_2^2
=\frac{\gamma^2}{K}\one_K^\top(AA^\top)^{-1}\one_K
\le4\gamma^2=\frac14.
\]
Thus $x_A$ is feasible and has support-function value $-1/(4\sqrt K)$. On every row array, $h_A^\star\le h_A(0)=0$. Averaging the event bound proves the expectation bound without any independence assumption.
\end{proof}

\section{The adaptive selected-coordinate posterior}\label{sec:posterior}

The hard distribution is fixed. All conditioning below is analytical: it does not change the objective or selection after interaction. We first state the external estimate and then verify that every posterior to which it is applied lies in its precise class.

\subsection{A one-step selection theorem}

\begin{definition}[Strong log-concavity on an affine support]\label{def:slc}
A probability measure $\mu$ on a nonempty affine subspace $H\subset\R^d$ is one-strongly log-concave on $H$ if
\begin{equation}\label{eq:slc-density}
\mu(du)=Z^{-1}e^{-V(u)}\sigma_H(du),\qquad 0<Z<\infty,
\end{equation}
where $\sigma_H$ is intrinsic Lebesgue measure, $V$ is proper and lower semicontinuous, and $u\mapsto V(u)-\norm{u}_2^2/2$ is convex on $H$. For a zero-dimensional $H$, $\sigma_H$ is unit mass at its single point.
\end{definition}

\begin{theorem}[External one-step selection inequality]\label{thm:external-selection}
\cite[Theorem~14]{ZZQL2026}.
Let $u_1,\ldots,u_k$ be independent random vectors, each with a law satisfying \cref{def:slc} on a possibly different affine support. Let $v\in\R^d$ have unit norm, and define
\[
Y=\max_{j\in[k]}\ip{u_j}{v},\qquad
J=\min\argmax_{j\in[k]}\ip{u_j}{v},\qquad S=\sum_{j=1}^k u_j.
\]
There is a universal constant $\Csel\ge1$ such that
\begin{equation}\label{eq:external-selection}
\E\norm{\E[S\mid J,Y]-\E S}_2^2\le\Csel\log(ek).
\end{equation}
The theorem includes $k=1$, constant projected scores, tied atoms, nonsmooth convex truncations, and zero-dimensional affine supports.
\end{theorem}

\begin{remark}[Affine-support formulation]
Equation~\eqref{eq:external-selection} is used as an external theorem, not proved here. Its general affine-support statement is necessary: after exact observations the posterior need not be full-dimensional or Gaussian without truncation. The one-row case is an estimate for the energy of a conditional mean after an exact linear observation. A covariance bound alone is not substituted for it.
\end{remark}

We also use parameterized disintegration on standard Borel spaces in the form of \cite[Lemma~10]{ZZQL2026}: a Borel family of probability laws and a jointly Borel observation map admit jointly Borel regular conditional kernels. For each parameter, the kernel is supported on the observation fiber almost surely. This permits conditioning on exact real values without treating null-level sets as positive-probability events.

\subsection{The augmented transcript}

Let the policy seed be $U$, independent of the rows. Define $W_0=U$. At time $t$, the policy chooses $(X_t,I_t)$ using its actual history. If $X_t\ne0$, define
\begin{equation}\label{eq:augmented-observation}
v_t=\frac{X_t}{\norm{X_t}_2},\qquad
J_t=\min\argmax_j\ip{b_j}{v_t},\qquad
M_t=\max_j\ip{b_j}{v_t},\qquad
C_t=(b_{J_t})_{I_t}.
\end{equation}
If $X_t=0$, define $(J_t,M_t)=(1,0)$ and $C_t=(b_1)_{I_t}$. Set
\begin{equation}\label{eq:W}
W_t=(W_{t-1},X_t,I_t,J_t,M_t,C_t),\qquad \cW_t=\sigma(W_t).
\end{equation}
The dummy value $M_t=0$ in the zero-query branch is deterministic and is not a hidden maximum in another direction.

The actual oracle response is
\begin{equation}\label{eq:response-equivalence}
[G_A(X_t)]_{I_t}
=\frac{C_t/\sqrt d+4r_{I_t}(X_t)}6.
\end{equation}
Conversely, $C_t=\sqrt d\bigl(6[G_A(X_t)]_{I_t}-4r_{I_t}(X_t)\bigr)$. Thus the coordinate $C_t$ is exactly equivalent to the actual scalar response, given the query. The actual history and output can be reconstructed from $W_t$. In particular, $(X_t,I_t)$ is $\cW_{t-1}$-measurable.

The label and score in \eqref{eq:augmented-observation} are not free oracle observations. Their conditional-mean energy cost will be charged explicitly, separately from the coordinate cost.

\begin{lemma}[Product posterior with convex affine slices]\label{lem:posterior}
For every fixed finite cap and admissible randomized policy, the augmented transcript is jointly Borel in $(A,U)$. For almost every completed transcript $W_t$, the conditional law of the rows factors into independent measures
\begin{equation}\label{eq:posterior-density}
\mu_{j,t}(db)=Z_{j,t}^{-1}e^{-\norm{b}_2^2/2}
\one_{C_{j,t}}(b)\,\sigma_{H_{j,t}}(db),\qquad 0<Z_{j,t}<\infty,
\end{equation}
where $H_{j,t}$ is a nonempty affine subspace and $C_{j,t}\subseteq H_{j,t}$ is a nonempty closed convex set. The conditional kernels can be chosen jointly Borel in the transcript. The same product form holds after revealing $(J_t,M_t)$ but before revealing $C_t$.
\end{lemma}

\begin{proof}
The initial law is a product of densities of the form \eqref{eq:posterior-density}, with $H_{j,0}=\R^d$ and $C_{j,0}=\{b:\norm{b}_2\le2\sqrt d\}$. Independence of $U$ leaves this law unchanged after conditioning on $W_0=U$. Joint Borelness follows inductively from \cref{lem:hard-borel} and the finite least-index maximum map.

Condition on a completed past for which the assertion holds. The query and requested coordinate are then fixed. If $X_t=0$, the label and score reveal no new information. Observing $C_t$ slices only row 1 by $\ip{b_1}{e_{I_t}}=C_t$, preserving product structure.

Suppose $X_t\ne0$. Given $(J_t,M_t)=(j,m)$, row $j$ is restricted by the affine equality
\begin{equation}\label{eq:winner-slice}
\ip{b_j}{v_t}=m,
\end{equation}
and the losing rows satisfy
\begin{equation}\label{eq:loser-truncation}
\ip{b_\ell}{v_t}<m\quad(\ell<j),\qquad
\ip{b_\ell}{v_t}\le m\quad(\ell>j).
\end{equation}
These are separate rowwise constraints. To justify conditional factorization also in the presence of atoms, let $\nu_j$ be the current law of $\ip{b_j}{v_t}$, and let $Q_j(m,db)$ be a regular conditional law of row $j$ given this score. Let $p_\ell^{(j)}(m)$ be the probability of the appropriate strict or weak event in \eqref{eq:loser-truncation}. Fubini's theorem gives the observation measure
\begin{equation}\label{eq:max-observation-measure}
\PP\{J_t=j,M_t\in dm\mid\cW_{t-1}\}
=\nu_j(dm)\prod_{\ell\ne j}p_\ell^{(j)}(m).
\end{equation}
Against this measure, the posterior is the product of $Q_j(m,db_j)$ and the normalized loser truncations. This identity follows first for products of row-test functions and then for all bounded Borel tests by a monotone-class argument. Normalizing probabilities that vanish occur on a null set for the observation measure. This proves independence after the label-and-score reveal.

The coordinate observation next imposes
\begin{equation}\label{eq:coordinate-slice}
\ip{b_j}{e_{I_t}}=C_t
\end{equation}
on the winner only. Applying rowwise disintegration again preserves the product law.

It remains to identify its geometry. A linear functional nonconstant on an affine support has level hyperplanes of zero intrinsic measure. Its almost-everywhere conditional law is the corresponding intrinsic affine slice, by the linear coarea formula. If the functional is constant on that support, its observation law is a point mass and no effective slice is added. Similarly, for a nonconstant loser score the strict and closed halfspaces differ by a null hyperplane. For a constant loser score the strict event is either the whole support or empty; empty branches have zero observation probability. Thus all effective truncations can be taken closed and convex on reached branches. The old supports acquire only affine equalities, and the old convex sets acquire only halfspaces or affine slices. Redundant equalities and zero-dimensional supports cause no exception. This gives \eqref{eq:posterior-density}.

Parameterized disintegration supplies jointly Borel conditional kernels. The fiberwise Gaussian-density identification follows in intrinsic coordinates, and the normalization is positive on almost every reached fiber. The additional measurable-coordinate details are given in \cref{app:kernels}. A finite intersection of the full-measure good-history sets suffices for the fixed cap. The induction is complete.
\end{proof}

\begin{corollary}[The external theorem remains applicable]\label{cor:posterior-slc}
Every row posterior in \cref{lem:posterior} satisfies \cref{def:slc}. Hence \cref{thm:external-selection} can be applied conditionally after every completed history and after every intermediate label-and-score reveal.
\end{corollary}
\begin{proof}
On $H_{j,t}$ its potential is $V(b)=\norm{b}_2^2/2+\ind_{C_{j,t}}(b)$, up to an additive normalizing constant. The convex indicator is proper, lower semicontinuous, and convex. A point-mass posterior is the zero-dimensional case of the same definition.
\end{proof}

\section{Posterior energy, expected risk, and the lower bound}\label{sec:lower}

\subsection{Charging both stages of an observation}

Let
\begin{equation}\label{eq:aggregate}
S=\sum_{j=1}^K b_j,\qquad S_{\mathrm{sc}}=\frac S{\sqrt d}=\sum_{j=1}^K a_j,
\qquad \Cen=2\Csel.
\end{equation}

\begin{theorem}[Selected-coordinate posterior-mean energy]\label{thm:energy}
For every admissible randomized coordinate-oracle policy with pathwise cap $T$,
\begin{equation}\label{eq:energy}
\E\norm{\E[S\mid\cW_T]}_2^2\le\Cen T\log(eK),\qquad
\E\norm{\E[S_{\mathrm{sc}}\mid\cW_T]}_2^2
\le\frac{\Cen T\log(eK)}d.
\end{equation}
There is no sparsity restriction on the query points and no restriction to uniform coordinate sampling.
\end{theorem}
\begin{proof}
The row law is centrally symmetric, and the seed is independent, so $\E[S\mid\cW_0]=0$. Let $m_t=\E[S\mid\cW_t]$ and insert the intermediate sigma-field
\[
\cW_{t-1}^{+}=\cW_{t-1}\vee\sigma(J_t,M_t),\qquad
m_t^+=\E[S\mid\cW_{t-1}^{+}].
\]
All variables are square-integrable, since the rows are bounded. Conditional on a good past, $v_t$ is fixed and the rows are independent with the affine-support laws of \cref{lem:posterior}. If $X_t\ne0$, \cref{thm:external-selection} gives
\begin{equation}\label{eq:score-energy}
\E[\norm{m_t^+-m_{t-1}}_2^2\mid\cW_{t-1}]
\le\Csel\log(eK).
\end{equation}
If $X_t=0$, this increment is zero, so the same inequality holds.

Conditional on $\cW_{t-1}^{+}$, the rows again form a product. Only the known winner row changes upon observing $C_t$. Apply \cref{thm:external-selection} with $k=1$ to that row, in direction $e_{I_t}$. Its maximum observation is just the scalar coordinate itself. Thus
\begin{equation}\label{eq:coordinate-energy}
\E[\norm{m_t-m_t^+}_2^2\mid\cW_{t-1}^{+}]\le\Csel.
\end{equation}
This includes a coordinate already constant on the posterior support.

The two increments are orthogonal in $L^2$: the first is measurable at the intermediate sigma-field, while the conditional expectation of the second given that field is zero. Consequently,
\[
\E\norm{m_t-m_{t-1}}_2^2
\le\Csel(\log(eK)+1)\le\Cen\log(eK).
\]
The martingale increments for different $t$ are also orthogonal. Summing over $t$ proves the first inequality in \eqref{eq:energy}; dividing by $d$ proves the second. The seed has been included in the initial sigma-field throughout, so the argument already covers arbitrary independent standard Borel randomization.
\end{proof}

\begin{remark}[Why augmentation is legitimate]
We do not assert that a scalar coordinate reveals only a fixed number of bits. Instead, the actual response is simulated exactly from the richer transcript, and both extra label-and-score information and the actual coordinate are charged in \eqref{eq:score-energy}--\eqref{eq:coordinate-energy}. A lower bound obtained from this richer information applies to the original transcript. No unqueried full row is revealed.
\end{remark}

\subsection{Direct conversion to expected optimization error}

\begin{proposition}[A finite-budget Bayes-risk inequality]\label{prop:bayes}
Suppose $d\ge1024$ and $1\le K\le d/256$. For every admissible randomized $T$-query policy on the normalized box,
\begin{equation}\label{eq:bayes-curve}
\E_{A,U}[F_A(\widehat x)-\min_{D_d}F_A]
\ge\frac7{192\sqrt K}
-\frac1{6K}\sqrt{\frac{\Cen T\log(eK)}d}.
\end{equation}
In particular, define
\begin{equation}\label{eq:bayes-constants}
c_E=\frac{49}{4096\Cen},\qquad \beta=\frac7{384}.
\end{equation}
If $T\le c_EdK/\log(eK)$, then
\begin{equation}\label{eq:bayes-lower}
\E_{A,U}[F_A(\widehat x)-\min_{D_d}F_A]\ge\frac\beta{\sqrt K}.
\end{equation}
\end{proposition}
\begin{proof}
Put $V_T=P(\widehat x)$, where $P$ is \eqref{eq:radial-projection}. It is $\cW_T$-measurable, has norm at most one, and has no larger loss than the original output. A maximum is at least the average, so
\[
h_A(V_T)\ge\frac1K\ip{S_{\mathrm{sc}}}{V_T}.
\]
Write $\overline S_T=\E[S_{\mathrm{sc}}\mid\cW_T]$. By conditional expectation, Cauchy--Schwarz, and \cref{thm:energy},
\begin{align}
\E h_A(V_T)
&\ge\frac1K\E\ip{\overline S_T}{V_T}
\ge-\frac1K\E\norm{\overline S_T}_2\notag\\
&\ge-\frac1K\sqrt{\E\norm{\overline S_T}_2^2}
\ge-\frac1K\sqrt{\frac{\Cen T\log(eK)}d}.\label{eq:mean-output}
\end{align}
Combine this with \cref{lem:negative-optimum} and the exact loss identity on the unit ball:
\[
F_A(V_T)-\min_{D_d}F_A=\frac16(h_A(V_T)-h_A^\star).
\]
This proves \eqref{eq:bayes-curve}. Under the stated cap,
\[
\frac1{6K}\sqrt{\frac{\Cen T\log(eK)}d}
\le\frac7{384\sqrt K},
\]
which subtracts half the positive term in \eqref{eq:bayes-curve} and proves \eqref{eq:bayes-lower}.
\end{proof}

The direct expectation argument needs no tail bound for the posterior residual. In particular, it does not condition the whole analysis on the negative-optimum event. Such conditioning would change the hard prior and could destroy the row-product structure. The spectral event enters only through the unconditional expectation bound \eqref{eq:negative-optimum}. A separate probability-of-failure argument, using a residual estimate, is included in \cref{app:residual}.

\subsection{Uniform choice of rows and minimax conversion}

\begin{theorem}[Normalized lower bound with specified constants]\label{thm:normalized}
With $c_E$ and $\beta$ from \eqref{eq:bayes-constants}, set
\begin{equation}\label{eq:final-constants}
a=\left(\frac\beta2\right)^2=\frac{49}{589824},\qquad
q_0=256,\qquad
d_0=\max\left\{1024,\left\lceil\frac2a\right\rceil\right\},\qquad
c=\frac{a c_E}{2}.
\end{equation}
Then, for every integer $d\ge d_0$ and every real $q\ge q_0$,
\begin{equation}\label{eq:normalized-lower}
\Tbox(d,q)\ge c\frac{d\min\{q^2,d\}}{\log(e\min\{q^2,d\})}.
\end{equation}
\end{theorem}
\begin{proof}
Put $s=\min\{q^2,d\}$ and choose
\begin{equation}\label{eq:K-choice}
K=\lfloor as\rfloor.
\end{equation}
The constants ensure $as\ge2$: indeed, $aq_0^2=49/9>2$ and $ad\ge2$. Therefore
\begin{equation}\label{eq:K-rounding}
\frac{as}{2}\le K\le as\le ad<\frac d{256}.
\end{equation}
Fix any integer cap with $T\le cds/\log(es)$. Since $K\le s$ and $K\ge as/2$,
\[
T\le c_E\frac{ads}{2\log(es)}
\le c_E\frac{dK}{\log(eK)}.
\]
For every policy, \cref{prop:bayes} applies to the same seed-independent product law and yields
\[
\E_{A,U}[F_A(\widehat x)-\min F_A]
\ge\frac\beta{\sqrt K}
\ge\frac\beta{\sqrt a\sqrt s}
=\frac2{\sqrt s}\ge\frac2q>\frac1q.
\]
The terminal loss is a jointly Borel nonnegative function of $(A,U)$ by \cref{lem:hard-borel}. Tonelli's theorem therefore permits the two integrations to be exchanged. If the policy had seed-expected error at most $1/q$ for every fixed pair in the support, averaging that inequality over $A$ would contradict the displayed strict bound. Hence there exists a deterministic complete pair $(F_A,G_A)$, fixed before the seed realization, on which the policy's expected error exceeds $1/q$.

This argument applies to every policy with the stated cap. No minimax interchange over an abstract function space is required. It is simply integration against a fixed finite-dimensional hard distribution. Legal padding treats early stopping. Thus no such integer cap is feasible in \eqref{eq:tbox}, proving \eqref{eq:normalized-lower}.
\end{proof}

\begin{remark}[Uniformity across the transition]
The choice \eqref{eq:K-choice} covers every $q\ge q_0$: $K\asymp q^2$ when $q\le\sqrt d$, and $K\asymp d$ when $q\ge\sqrt d$. The rounding bound \eqref{eq:K-rounding} also holds at the transition. The thresholds above may be taken as $q_0=256$ and $d_0=24075$; they are conservative and are not optimized. The lower-bound coefficient $c$ depends on the universal constant in the cited selection theorem, whose numerical value is not needed.
\end{remark}

\subsection{Scaling to general problem parameters}

\begin{lemma}[Exact coordinate-oracle scaling]\label{lem:scaling}
For $G,R_B>0$, each normalized hard pair defines a physical pair on $R_B[-1,1]^d$ by
\begin{equation}\label{eq:physical-scaling}
\mathcal F_A(x)=GR_B F_A(x/R_B),\qquad
\mathcal G_A(x)=G G_A(x/R_B).
\end{equation}
It belongs to $\cI(R_B[-1,1]^d,0;G,R_B)$. Feasible policies, outputs, scalar responses, and expected gaps correspond exactly under this scaling, with one physical scalar call per normalized scalar call. Consequently,
\begin{equation}\label{eq:physical-lower}
\Tcoord(d,G,R_B,GR_B/q)\ge\Tbox(d,q).
\end{equation}
\end{lemma}
\begin{proof}
Substitute $y=x/R_B$ in the normalized subgradient inequality to obtain, for physical $x,z$,
\[
\mathcal F_A(z)\ge\mathcal F_A(x)+\ip{\mathcal G_A(x)}{z-x}.
\]
The selection norm is at most $G$, and the objective is $G$-Lipschitz. A normalized optimizer of norm at most one scales to an optimizer $x^\star$ of norm at most $R_B$, so the radius promise holds. Borelness and coherence are preserved by the continuous scaling maps.

The same formulas also scale every pair in $\cI([-1,1]^d,0;1,1)$, not just the hard family: its promised optimizer has norm at most $\sqrt2$ and scales to one of norm at most $\sqrt2R_B$, exactly as required in \eqref{eq:radius}.

A normalized query $(y,i)$ corresponds to $(R_By,i)$; its response is multiplied by $G$. Division by the known positive $G$ recovers the normalized response exactly. The reverse simulation is equally direct. Output scaling gives
\[
\mathcal F_A(R_By)-\min\mathcal F_A
=GR_B(F_A(y)-\min F_A).
\]
Thus tolerance $1/q$ becomes $GR_B/q$, and neither the number of scalar calls nor the pathwise cap changes. The scaled box is one public geometry in \eqref{eq:minimax}. A policy that worked uniformly over all public geometries at a smaller cap would, on this box, pull back to a normalized policy contradicting \eqref{eq:tbox}. This proves \eqref{eq:physical-lower}.
\end{proof}

\begin{proof}[Proof of \cref{thm:main}]
The lower bound is \cref{thm:normalized,lem:scaling}. The first upper bound follows from \cref{thm:rcs}; the second follows from \cref{thm:centroid,lem:cp-borel}. Selecting the smaller of the two certified integer budgets uses only the public parameters. The deterministic cutting-plane guarantee is also an expected-error guarantee. The simplified logarithmic expression follows from \cref{cor:cp-log}.
\end{proof}

\section{Scope and remaining questions}\label{sec:scope}

\subsection{Polynomial dependence and logarithmic gaps}

For the parameter range of the lower bound,
\begin{equation}\label{eq:frontier-final}
c\frac{d\min\{q^2,d\}}{\log(e\min\{q^2,d\})}
\le\Tcoord(d,G,R_B,\eps)
\le C\min\{dq^2,d^2\log(edq)\}.
\end{equation}
The same form holds for $\Tupd$ by \cref{cor:update}. One may summarize either statement as $\widetilde\Theta(d\min\{q^2,d\})$ when the suppressed factors are allowed to depend logarithmically on both $d$ and $q$. The explicit bounds in \eqref{eq:frontier-final} are more informative at very high accuracy. They leave a logarithmic gap in the regime $q\le\sqrt d$ and do not establish a lower dependence on $\log q$ when $q\ge\sqrt d$.

\subsection{The information returned by the oracle}

The oracle returns a coordinate of a selected vector subgradient. This differs from exact coordinate minimization, access to an entire univariate restriction, or a one-sided directional derivative. For example, for $f(x_1,x_2)=\max\{x_1,x_2\}$, the vector $e_1$ is a valid selected subgradient at zero, but the one-sided derivative in direction $-e_1$ is zero. Thus the selected component does not by itself specify both directional derivatives along that coordinate.

Coherence is also essential to the upper bound. Querying all coordinates at a fixed point reconstructs one valid subgradient because the selection is fixed. If each scalar answer could instead be chosen from a different subgradient, the assembled vector might fail the subgradient inequality. Our lower bound already holds for the canonical selection \eqref{eq:hard-pair}, fixed before the seed is realized.

\subsection{Coordinate access, coordinatewise motion, and update rules}

A coordinate method can have dense iterates; changing one coordinate at a time does not impose a sparse-query model. The lower bound covers arbitrary points in the box and therefore applies to all coordinate-update policies in \cref{def:update}.

Conversely, coordinatewise motion does not prohibit global computation. \Cref{lem:routing,cor:routed-centroid} show that centroid localization can be implemented under this motion restriction with at most twice as many scalar queries. Hence the broad coordinate-update class has the same polynomial complexity in both regimes. A prescribed stepsize rule, subgradient recursion, or memory restriction can define a narrower class. In particular, last-iterate guarantees for a prescribed subgradient recursion do not follow from the routing construction. No such additional restriction is imposed in \cref{def:update}.

\subsection{Geometry and the optimizer-distance promise}

The parameter $R_B$ controls the distance from $x_1$ to one optimizer through \eqref{eq:radius}; it does not bound the diameter of the whole domain. The scaled hard box has diameter $2R_B\sqrt d$, while the radial penalty places all its optimizers within distance $R_B$ of the initial point. The upper bound localizes to $D\cap\overline B_2(x_1,\sqrt2R_B)$ and is therefore valid even when $D$ is arbitrarily wider. Coordinatewise routing may leave this localized set, but always remains in the legal domain $D$.

The minimax definition takes the worst case over known product geometries. It does not claim the same lower bound for each geometry: singleton domains are trivial. A class whose entire domain has diameter at most a fixed multiple of $R_B$ would require a separate lower-bound analysis.

\subsection{Limitations and dependencies}

The upper bounds hold in every dimension, whereas the lower bound requires universal thresholds on $d$ and $q$. The present argument does not give sharp lower bounds for small fixed dimensions or determine the optimal logarithmic factors. A model that limits only the expected number of queries also requires a separate analysis. The present results use a deterministic cap on each execution.

The lower-bound proof relies on the one-step selection theorem of Zhang et al.\ \cite[Theorem~14]{ZZQL2026}. We state that theorem precisely and verify its hypotheses under each adaptive conditioning step. The selected-coordinate extension, radial localization, expected-risk conversion, and scaling arguments are proved here. An independent proof of the cited analytic theorem is outside the scope of this paper. The optional probability result in \cref{app:residual} additionally uses a standard strong-log-concavity moment bound, stated explicitly there.

\section{Conclusion}

Exact coordinate access has an intrinsic dimension-dependent cost in convex Lipschitz optimization. Even with dense adaptive queries, arbitrary independent randomization, and a fixed coherent subgradient selection, the worst-instance complexity has polynomial scale $d\min\{(GR_B/\eps)^2,d\}$. The lower bound follows by charging the posterior-mean energy of both the active-row selection and the subsequent coordinate observation. Projected coordinate subgradient updates and value-free centroid localization supply the complementary upper bounds, and coordinatewise routing extends the latter to policies with coordinatewise motion. Further work can refine the logarithmic factors and analyze more restrictive algorithm classes; the lower bounds proved here apply only in the stated dimension and accuracy range.

\appendix
\section{Measurable kernels and exact affine conditioning}\label{app:kernels}

This appendix expands the null-fiber and joint-measurability points used in \cref{lem:posterior}. All assertions concern finite-dimensional row variables and a standard Borel transcript parameter.

\subsection{Parameterized disintegration}

The standard disintegration statement used here is as follows. Let $\Theta,X,Y$ be standard Borel spaces, $K(\theta,dx)$ a probability kernel, and $T:\Theta\times X\to Y$ Borel. If $K_T(\theta,dy)$ is its observation image, there is a jointly Borel probability kernel $Q(\theta,y,dx)$ such that, for every nonnegative Borel $h$,
\begin{equation}\label{eq:disintegration}
\begin{split}
&\int h(\theta,T(\theta,x),x)K(\theta,dx)\\
&\qquad=\iint h(\theta,y,x)Q(\theta,y,dx)K_T(\theta,dy).
\end{split}
\end{equation}
For each parameter, $Q$ is a regular conditional law and is supported on the observation fiber almost surely. This is the form recorded in \cite[Lemma~10]{ZZQL2026}. The identity, rather than division by the probability of a single level set, defines conditioning on exact real observations.

\subsection{An explicit one-row slice}

Consider a posterior row with affine support $H=b_0+L$, where $b_0$ is the minimum-norm point of $H$, and let $Q\in\R^{d\times r}$ have orthonormal columns spanning $L$. Its density in $z\in\R^r$ is proportional to
\[
\exp\left(-\frac12\norm{b_0+Qz}_2^2\right)\one_C(b_0+Qz).
\]
Observe $\ip{b}{v}=m$ and put $a=Q^\top v$. If $a=0$, the score is the constant $\ip{b_0}{v}$ on all of $H$; at that observed value the conditional row law is unchanged.

If $a\ne0$, choose $n_a=a/\norm a_2$ and an orthonormal matrix $R_a$ whose columns span $a^\perp\subset\R^r$. Write
\[
z=n_a\,t+R_a w,
\qquad
t=\frac{m-\ip{b_0}{v}}{\norm a_2}.
\]
Conditionally on $m$, the row lies on
\begin{equation}\label{eq:explicit-fiber}
b=b_0+Qn_a\frac{m-\ip{b_0}{v}}{\norm a_2}+QR_a w,
\qquad w\in\R^{r-1},
\end{equation}
with density proportional to $e^{-\norm b_2^2/2}\one_C(b)$ relative to this intrinsic coordinate. The coarea factor $1/\norm a_2$ is constant in $w$ and cancels upon normalization. The new feasible set is the closed convex set $C\cap\{b:\ip{b}{v}=m\}$. The normalizer is positive for almost every observed $m$. If $r=1$, the coordinate $w$ is absent and the slice is a point mass.

This shows directly why equality observations preserve the required strong log-concavity. It also covers repeated observations: a previously imposed equality makes the next corresponding functional constant, and no new dimension is removed.

\subsection{Joint Borel coordinate choices}

At finite time each support is specified by finitely many affine equations with coefficients Borel in the transcript. The rank of its coefficient matrix is Borel, because rank strata can be described using finitely many minors. On each consistent rank stratum, Gaussian elimination using the lexicographically first admissible pivot pattern gives a Borel particular solution. Orthogonal projection or the corresponding finite pseudoinverse formula gives the minimum-norm solution $b_0$.

A Borel orthonormal basis of the nullspace can be obtained by projecting the standard basis vectors, taking the first linearly independent ones in lexicographic order, and applying Gram--Schmidt with normalization only on the nonzero branches. The same construction gives $R_a$ in \eqref{eq:explicit-fiber}. There are finitely many possible ranks and pivot patterns for a fixed cap. The coefficients in the intrinsic density are therefore Borel on every branch.

Integrals of nonnegative jointly Borel functions over the fixed Euclidean coordinates are Borel in the parameters. Hence the slice normalizers, moment integrals, and conditional kernels admit Borel versions on the positive-normalizer branches. On inconsistent affine systems or zero-normalizer observations, choose an arbitrary fixed probability measure to complete the kernel definition. Such branches are null for the actual observation law and do not alter \eqref{eq:disintegration}.

\subsection{Maximum observations and ties}

Formula~\eqref{eq:max-observation-measure} treats continuous and atomic score laws in the same way. For a possible winner $j$, the constraints on earlier labels are strict, whereas those on later labels are weak. If a loser functional is nonconstant on its support, the boundary hyperplane is intrinsically null, so closing a strict halfspace does not change that row's truncated law. If the functional is constant, the strict event is either the whole support or empty; the empty case carries zero mass in \eqref{eq:max-observation-measure}.

This distinction is necessary after adaptive conditioning: even if the initial row law is absolutely continuous in $\R^d$, later scores may be deterministic or tied. The least-index rule is never perturbed or selected anew. On each reached branch, the posterior remains a product of closed convex affine slices under the original fixed tie convention.

\subsection{Why arbitrary seeds do not change the argument}

One may include $U$ as the first coordinate of every transcript, as done in the main proof, or first fix a seed and then integrate. In the latter formulation every deterministic section has the same energy and risk constants, and the row law is independent of that section. Joint Borelness of the loss and Tonelli's theorem justify integration over any standard Borel seed law. Neither approach allows the hard rows or selection to depend on the realized seed.

\section{A probability-of-failure lower bound}\label{app:residual}

The main expected-error proof does not need this appendix. We record a second conversion, using a posterior-residual estimate, to distinguish probability and expectation statements explicitly.

\subsection{The additional external moment bound}

A one-strongly log-concave law on $H=b_0+L$ satisfies
\begin{equation}\label{eq:external-subgaussian}
\E\exp\{\lambda\ip{b-\E b}{v}\}
\le\exp\left(\frac{\lambda^2\norm{P_Lv}_2^2}{2}\right)
\qquad(\lambda\in\R,\ v\in\R^d).
\end{equation}
This is \cite[Lemma~37]{ZZQL2026}, including zero-dimensional affine supports. It is an additional external input only for this appendix, not for \cref{prop:bayes,thm:normalized}.

\begin{lemma}[Residual in a transcript-dependent direction]\label{lem:residual}
Let
\[
\overline S_T=\E[S_{\mathrm{sc}}\mid\cW_T],\qquad
Z_T=S_{\mathrm{sc}}-\overline S_T.
\]
For any $\cW_T$-measurable $V_T\in\ball$ and every $u>0$,
\begin{equation}\label{eq:residual-tail}
\PP\{\ip{Z_T}{V_T}\le-u\mid\cW_T\}
\le\exp\left(-\frac{du^2}{2K}\right)
\quad\text{almost surely}.
\end{equation}
\end{lemma}
\begin{proof}
Condition on a good completed transcript. The rows are independent and satisfy \eqref{eq:external-subgaussian}. The direction $V_T$ is now fixed and has norm at most one. Therefore
\begin{align*}
\E[e^{\lambda\ip{Z_T}{V_T}}\mid\cW_T]
&=\prod_{j=1}^K\E\left[
\exp\left\{\frac\lambda{\sqrt d}\ip{b_j-\E[b_j\mid\cW_T]}{V_T}\right\}
\middle|\cW_T\right]\\
&\le\exp\left(\frac{\lambda^2K}{2d}\right).
\end{align*}
Use conditional Chernoff's inequality with $-\lambda$, and minimize at $\lambda=du/K$.
\end{proof}

\begin{proposition}[A fixed-prior probability lower bound]\label{prop:probability}
Let $\gamma=1/4$, $\eta_0=\gamma/16=1/64$, and
\[
c_P=\frac{\eta_0^2}{16\Cen}.
\]
Suppose
\[
d\ge\max\left\{1024,\left\lceil\frac{2\log16}{\eta_0^2}\right\rceil\right\},
\qquad 1\le K\le d/256,
\qquad T\le c_P\frac{dK}{\log(eK)}.
\]
Then, under the same fixed product prior and any independent policy seed,
\begin{equation}\label{eq:prob-failure}
\PP_{A,U}\left\{F_A(\widehat x)-\min F_A\ge\frac7{192\sqrt K}\right\}
\ge\frac34.
\end{equation}
Consequently its Bayes expected error is at least $7/(256\sqrt K)$.
\end{proposition}
\begin{proof}
Put $V_T=P(\widehat x)$. Define
\[
E_{\mathrm{mean}}=\{\norm{\overline S_T}_2\le\eta_0\sqrt K\},
\qquad
E_{\mathrm{res}}=\{\ip{Z_T}{V_T}\ge-\eta_0\sqrt K\}.
\]
By \cref{thm:energy}, Markov's inequality, and the cap,
\[
\PP(E_{\mathrm{mean}}^c)
\le\frac{\Cen T\log(eK)}{d\eta_0^2K}\le\frac1{16}.
\]
By \cref{lem:residual} and the dimension assumption,
\[
\PP(E_{\mathrm{res}}^c)\le e^{-d\eta_0^2/2}\le\frac1{16}.
\]
On the intersection,
\[
h_A(V_T)\ge\frac1K\ip{S_{\mathrm{sc}}}{V_T}
\ge-\frac{2\eta_0}{\sqrt K}
=-\frac\gamma{8\sqrt K}.
\]
The negative-optimum event of \cref{lem:negative-optimum} has probability at least $7/8$. On its intersection with the preceding two events,
\[
F_A(V_T)-\min F_A
\ge\frac16\left(\frac\gamma{\sqrt K}-\frac\gamma{8\sqrt K}\right)
=\frac7{192\sqrt K}.
\]
A union bound, not independence, gives intersection probability at least $1-1/8-1/16-1/16=3/4$. The radial projection cannot increase loss. This proves \eqref{eq:prob-failure}. Multiplying the loss threshold by $3/4$ and using nonnegativity gives the expectation bound.
\end{proof}

\begin{remark}
The larger expected-loss constant in this appendix comes with a smaller allowed cap and an additional moment-bound dependency. It does not replace the simpler main proof. In both versions, the hard distribution is chosen before the seed, and the final worst-instance statement follows only after the appropriate averaging argument.
\end{remark}

\section{Exact finite-transcript interpolation}\label{app:interpolation}

This appendix characterizes when a finite coordinate transcript is consistent with a convex Lipschitz objective, a fixed bounded subgradient selection, and a proposed optimizer. The max-affine construction and cycle criterion are finite-data forms of classical convex interpolation and cyclic monotonicity; see Rockafellar\ \cite[Theorem~1]{Rockafellar1966}. Here they are formulated with partially observed coordinates, a common selection at repeated query points, and an explicit optimizer-distance constraint. The results are auxiliary to the lower bound and also identify a nonconvexity that arises when the optimizer location is unknown.

\subsection{Fixed-optimizer fibers}

Let $D\subset\R^d$ be a nonempty compact convex set, and fix $z\in D$. Let $x^1,\ldots,x^m\in D$ be distinct queried points. At each point prescribe a set of coordinate observations $v_i^s=y_{s,i}$. Repeated locations are merged, and repeated observations of the same subgradient coordinate must agree.

\begin{proposition}[Exact fixed-optimizer interpolation]\label{prop:interpolation}
The following are equivalent.
\begin{enumerate}[label=(\roman*)]
\item There is a finite convex $G$-Lipschitz function $F:D\to\R$ with a fixed Borel selection $h$ satisfying \eqref{eq:coherence}, realizing the prescribed coordinate observations, and having $z\in\argmin_DF$.
\item There are scalars $a_*,a_1,\ldots,a_m$ and vectors $v^1,\ldots,v^m$ such that
\begin{align}
\norm{v^s}_2&\le G,\qquad v_i^s=y_{s,i}\text{ on observed coordinates},\label{eq:interp-norm}\\
a_r&\ge a_s+\ip{v^s}{x^r-x^s}\qquad(r,s\in[m]),\label{eq:interp-pair}\\
a_*&\le a_s,\qquad a_s+\ip{v^s}{z-x^s}\le a_*
\qquad(s\in[m]).\label{eq:interp-opt}
\end{align}
\end{enumerate}
For fixed $z$, the feasible set in these variables is closed and convex.
\end{proposition}
\begin{proof}
For necessity, take $a_s=F(x^s)$, $a_*=F(z)$, and $v^s=h(x^s)$. The norm bound, subgradient inequalities, and optimality give the three sets of constraints.

For sufficiency, define on all of $\R^d$
\begin{equation}\label{eq:maxaffine-completion}
F(x)=\max\left\{a_*,\ a_s+\ip{v^s}{x-x^s}:s\in[m]\right\}.
\end{equation}
It is finite and convex, and all slopes have norm at most $G$, so it is $G$-Lipschitz. At $x^r$, every affine piece is at most $a_r$ by \eqref{eq:interp-pair} and \eqref{eq:interp-opt}, while the $r$th piece equals $a_r$. Hence $F(x^r)=a_r$ and $v^r\in\partial F(x^r)$. At $z$, all affine pieces are at most $a_*$, and the constant piece gives equality. Thus $F(z)=a_*$ is the global minimum.

At each queried point choose $h(x^s)=v^s$. Away from the finitely many queried points, choose the slope of the least-index active affine piece, assigning slope zero to the constant piece. This is a Borel selection, because its cells are determined by finitely many affine comparisons. Its finite patching at the queried points preserves Borelness. Every selected slope is active, so it is a valid subgradient with norm at most $G$. Restriction to $D$ proves (i). Finally, for fixed $z$ the constraints are linear inequalities and Euclidean ball constraints, proving closed convexity.
\end{proof}

\subsection{Cycle criterion}

Normalize $a_*=0$ and adjoin a formal optimizer node $(x^0,v^0)=(z,0)$. If $z$ is a queried point, retain this formal node separately: it represents the constant supporting plane, not necessarily the actual selected subgradient at $z$.

\begin{proposition}[No-positive-cycle characterization]\label{prop:cycles}
A completion at fixed $z$ exists if and only if there are vectors $v^1,\ldots,v^m$ satisfying the norm and observed-coordinate constraints for which every directed cycle $s_0,s_1,\ldots,s_p=s_0$ on $\{0,\ldots,m\}$ satisfies
\begin{equation}\label{eq:cycles}
\sum_{r=0}^{p-1}\ip{v^{s_r}}{x^{s_{r+1}}-x^{s_r}}\le0.
\end{equation}
\end{proposition}
\begin{proof}
Set $c_{sr}=\ip{v^s}{x^r-x^s}$. The interpolation inequalities, including node 0, are $a_r\ge a_s+c_{sr}$. Summing around a cycle proves necessity.

For sufficiency, define $a_r$ as the maximum weight of a directed path from node 0 to node $r$. Under \eqref{eq:cycles}, removing a cycle cannot decrease path weight, so the maximum is attained among simple paths and is finite. The empty path gives $a_0=0$; no cycle returning to 0 has positive weight. Appending the edge $s\to r$ to a maximizing path and deleting repeated cycles shows $a_r\ge a_s+c_{sr}$. These are exactly \eqref{eq:interp-pair} and \eqref{eq:interp-opt} with $a_*=0$. Apply \cref{prop:interpolation}.
\end{proof}

\subsection{The exact worst gap consistent with a transcript}

\begin{proposition}[Finite-dimensional transcript risk]\label{prop:transcript-risk}
Fix a candidate output $y\in D$ and a proposed optimizer $z$. Adjoin $y$ as an additional node with latent vector $v^y$ and value $a_y$, or use its existing node if it was queried. Retain the formal optimizer node $(z,0)$, fix $a_0=0$, and impose all pairwise interpolation inequalities and norm constraints. If the resulting fiber is feasible, then
\begin{equation}\label{eq:exact-transcript-risk}
\max a_y
=\sup\{F(y)-F(z):(F,h)\text{ realizes the transcript},\ z\in\argmin_DF\}.
\end{equation}
The maximum on the left exists. Taking the supremum over feasible $z\in D\cap\overline B_2(x_1,\sqrt2R_B)$ gives the exact worst output gap consistent with the transcript and radius promise.
\end{proposition}
\begin{proof}
Every consistent pair supplies feasible variables with $a_s=F(x^s)-F(z)$ and $v^s=h(x^s)$, so its gap is bounded by the left side. Conversely, every feasible choice is realized by \eqref{eq:maxaffine-completion}, including the output node, with gap exactly $a_y$. The feasible set is compact: all vectors lie in closed norm-$G$ balls, while the inequalities between node 0 and any other node imply
\[
0\le a_s\le G\norm{x^s-z}_2.
\]
Thus a maximizing choice exists. The final assertion follows because an admissible consistent pair has at least one optimizer in the displayed promise set, and every feasible fixed-$z$ completion meets that promise.
\end{proof}

\begin{proposition}[The optimizer union is not generally convex]\label{prop:nonconvex-fibers}
The union of fixed-optimizer fibers need not be convex in the joint variables consisting of optimizer location and latent subgradient vectors. It cannot, in general, be represented by a convex lift whose linear projection retains those variables exactly.
\end{proposition}
\begin{proof}
Let $G,R>0$, $D=\{0\}\times[0,R]$, and query the first subgradient coordinate at $x=(0,0)$, observing zero. For a latent vector $v=(0,w)$ and a proposed optimizer $z=(0,t)$, \eqref{eq:interp-opt} implies
\[
a_1+wt\le a_*\le a_1,\qquad\text{hence }wt\le0.
\]
The pair $(w,t)=(G,0)$ is feasible, realized by $F(0,y)=Gy$ with its indicated subgradient selection. The pair $(w,t)=(0,R)$ is feasible for the constant zero function. Their midpoint $(G/2,R/2)$ violates $wt\le0$. Thus the joint feasible projection is nonconvex. Since a linear image of a convex set is convex, no exact convex lift with that projection exists. To include the radius promise, choose $R\le\sqrt2R_B$ and $x_1=0$.
\end{proof}

\begin{remark}[Consistency is not a seed-independent lower bound]
These interpolation results describe possible completions of a finite transcript. Choosing a different completion after observing the seed would not establish a randomized lower bound in \cref{def:complexity}. The main proof instead fixes the distribution of complete row arrays before the seed is realized. Also, nonconvexity of the optimizer union does not rule out all algorithms or relaxations; it rules out only the specified exact convex representation.
\end{remark}

\begin{thebibliography}{99}

\bibitem{ZZQL2026}
Haihan Zhang, Chenheng Zhang, Zhiquan Qi, and Zhouchen Lin.
\newblock Near-optimal lower bounds for randomized algorithms in exact value zeroth-order convex optimization.
\newblock \emph{arXiv preprint arXiv:2607.16558v2}, 2026.
\newblock \href{https://arxiv.org/abs/2607.16558v2}{arXiv:2607.16558v2}.
\newblock Version of July 30, 2026.

\bibitem{Nesterov2012}
Yurii Nesterov.
\newblock Efficiency of coordinate descent methods on huge-scale optimization problems.
\newblock \emph{SIAM Journal on Optimization}, 22(2):341--362, 2012.
\newblock \href{https://doi.org/10.1137/100802001}{doi:10.1137/100802001}.

\bibitem{DangLan2015}
Cong D. Dang and Guanghui Lan.
\newblock Stochastic block mirror descent methods for nonsmooth and stochastic optimization.
\newblock \emph{SIAM Journal on Optimization}, 25(2):856--881, 2015.
\newblock \href{https://doi.org/10.1137/130936361}{doi:10.1137/130936361}.

\bibitem{DengLanRangarajan2015}
Qi Deng, Guanghui Lan, and Anand Rangarajan.
\newblock Randomized block subgradient methods for convex nonsmooth and stochastic optimization.
\newblock \emph{arXiv preprint arXiv:1509.04609}, 2015.
\newblock \href{https://arxiv.org/abs/1509.04609}{arXiv:1509.04609}.

\bibitem{ZhaoEtAl2023}
Lei Zhao, Ding Chen, Daoli Zhu, and Xiao Li.
\newblock Randomized coordinate subgradient method for nonsmooth composite optimization.
\newblock \emph{arXiv preprint arXiv:2206.14981v3}, 2023.
\newblock \href{https://arxiv.org/abs/2206.14981v3}{arXiv:2206.14981v3}.

\bibitem{AGMSMGYN2024}
David Alonso-Guti\'errez, Francisco Mar\'in Sola, Javier Mart\'in Go\~ni, and Jes\'us Yepes Nicol\'as.
\newblock A general functional version of Gr\"unbaum's inequality.
\newblock \emph{Journal of Mathematical Analysis and Applications}, 544(1):129065, 2025.
\newblock \href{https://doi.org/10.1016/j.jmaa.2024.129065}{doi:10.1016/j.jmaa.2024.129065}.
\newblock Preprint: \href{https://arxiv.org/abs/2404.08319v1}{arXiv:2404.08319v1}.

\bibitem{Naszodi2017}
M\'arton Nasz\'odi.
\newblock Approximating a convex body by a polytope using the epsilon-net theorem.
\newblock \emph{Discrete \& Computational Geometry}, 61(3):686--693, 2019.
\newblock \href{https://doi.org/10.1007/s00454-018-9977-0}{doi:10.1007/s00454-018-9977-0}.

\bibitem{Rockafellar1966}
R. T. Rockafellar.
\newblock Characterization of the subdifferentials of convex functions.
\newblock \emph{Pacific Journal of Mathematics}, 17(3):497--510, 1966.
\newblock \href{https://doi.org/10.2140/pjm.1966.17.497}{doi:10.2140/pjm.1966.17.497}.

\end{thebibliography}
\end{document}

